\documentclass[11pt,a4paper]{article}
\usepackage[margin=20mm,top=20mm,bottom=20mm,headheight=15pt]{geometry}
\usepackage{fontspec}
\usepackage{amsmath,amssymb,mathtools}
\usepackage{unicode-math}
\setmainfont{Noto Serif CJK SC}[BoldFont={Noto Serif CJK SC Bold},Mapping=tex-text]
\setsansfont{Noto Sans CJK SC}[BoldFont={Noto Sans CJK SC Bold},Mapping=tex-text]
\XeTeXlinebreaklocale "zh"
\XeTeXlinebreakskip=0pt plus 1pt
\renewcommand{\contentsname}{目录}
\setmonofont{Latin Modern Mono}
\setmathfont{Latin Modern Math}
\usepackage{graphicx,xcolor}
\definecolor{Ink}{HTML}{20282D}
\definecolor{Muted}{HTML}{626B70}
\usepackage{longtable,booktabs,array,calc}
\usepackage{caption,float,placeins}
\renewcommand{\topfraction}{0.95}
\renewcommand{\bottomfraction}{0.95}
\renewcommand{\textfraction}{0.05}
\renewcommand{\floatpagefraction}{0.65}
\setlength{\intextsep}{8pt plus 2pt minus 2pt}
\setlength{\textfloatsep}{10pt plus 2pt minus 2pt}
\setlength{\floatsep}{10pt plus 2pt minus 2pt}
\usepackage{enumitem,fancyhdr,titlesec}
\usepackage{needspace}
\usepackage{microtype}
\usepackage{xurl}
\usepackage[unicode,hidelinks,breaklinks]{hyperref}
\hypersetup{pdftitle={补充材料：认知预测地图中方向与时间权重的机制识别},pdfauthor={Yongzhang
Ji},pdfsubject={SR mechanism identification, experimental design and joint measurement uncertainty},pdfkeywords={successor representation, cognitive maps, identifiability, qualitative properties, temporal direction}}
\setlength{\parindent}{0pt}
\setlength{\parskip}{5pt plus 1pt minus 1pt}
\linespread{1.12}
\lefthyphenmin=3
\righthyphenmin=3
\hyphenpenalty=1500
\exhyphenpenalty=1500
\setlength{\emergencystretch}{3em}
\setcounter{secnumdepth}{0}
\setcounter{tocdepth}{2}
\raggedbottom
\widowpenalty=10000
\clubpenalty=10000
\displaywidowpenalty=10000
\allowdisplaybreaks[1]
\predisplaypenalty=10000
\postdisplaypenalty=1000
\titleformat{\section}{\sffamily\bfseries\Large\color{Ink}}{}{0em}{}
\titleformat{\subsection}{\sffamily\bfseries\normalsize\color{Ink}}{}{0em}{}
\titleformat{\subsubsection}{\sffamily\bfseries\normalsize}{}{0em}{}
\titlespacing*{\section}{0pt}{17pt plus 3pt}{7pt}
\titlespacing*{\subsection}{0pt}{12pt plus 2pt}{5pt}
\titlespacing*{\subsubsection}{0pt}{10pt plus 2pt}{4pt}
\setlist{nosep,leftmargin=1.5em}
\setlist[enumerate]{itemsep=3pt}
\providecommand{\tightlist}{\setlength{\itemsep}{0pt}\setlength{\parskip}{0pt}}
\providecommand{\passthrough}[1]{#1}
\makeatletter
\def\maxwidth{\ifdim\Gin@nat@width>\linewidth\linewidth\else\Gin@nat@width\fi}
\def\maxheight{\ifdim\Gin@nat@height>0.56\textheight0.56\textheight\else\Gin@nat@height\fi}
\makeatother
\setkeys{Gin}{width=\maxwidth,height=\maxheight,keepaspectratio}
\pagestyle{fancy}
\fancyhf{}
\fancyhead[L]{\sffamily\footnotesize\color{Muted} 认知预测地图的机制识别}
\fancyhead[R]{\sffamily\footnotesize\color{Muted} 补充材料}
\fancyfoot[C]{\sffamily\footnotesize\color{Muted}\thepage}
\renewcommand{\headrulewidth}{0pt}
\renewcommand{\footrulewidth}{0pt}
\begin{document}
\thispagestyle{plain}
{\raggedright\hyphenpenalty=10000\sffamily\bfseries\color{Ink}\fontsize{20}{25}\selectfont 补充材料：认知预测地图中方向与时间权重的机制识别\par}
\vspace{8pt}
{\sffamily\fontsize{12}{15}\selectfont Yongzhang Ji\par}
\vspace{4pt}
{\sffamily\small Faculty of Psychology, Southwest University, Chongqing,
China\par}
\vspace{3pt}
{\sffamily\footnotesize\color{Muted} ORCID: \href{https://orcid.org/0009-0006-4111-2000}{0009-0006-4111-2000}\par}
\vspace{7pt}
{\sffamily\footnotesize\color{Muted} 补充材料\quad |\quad 2026-11-26\par}
\vspace{8pt}
本文件提供正文各结论的一般性证明、真实重复测量与已知变换校准的完整方法和结果，以及相关研究背景。理论结果
T01–T28 与正文相互索引；经验材料 E1–E11
汇集全部分析分支及其数据来源。图的矢量源文件、源数据和复现代码随文提供。

\begingroup
\fontsize{8.8}{10.3}\selectfont
\linespread{1.0}\selectfont
\setlength{\parskip}{0pt}
\tableofcontents
\endgroup

\clearpage

\section{A.
理论与统计证明}\label{a.-ux7406ux8bbaux4e0eux7edfux8ba1ux8bc1ux660e}

本附录按“观测对象—不可辨边界—实验解除—可观测读出—不确定性”组织证明。所有方向声明均相对于预先规定的状态标签与转移方向。各命题的编号
T01–T28
在正文、理论清单及复现文件中一致。这里区分三种不同目标：恢复任意完整算子；辨认受限模型族共同具有的性质；在该族内估计参数。三者需要的信息不相同。

\Needspace{7\baselineskip}
\subsection{A1.
模型、状态对比与观测层级}\label{a1.-ux6a21ux578bux72b6ux6001ux5bf9ux6bd4ux4e0eux89c2ux6d4bux5c42ux7ea7}

\Needspace{6\baselineskip}
\subsubsection{A1.1 符号}\label{a1.1-ux7b26ux53f7}

\par
\noindent\begin{minipage}{\linewidth}

设 \(P\in\mathbb R^{n\times n}\) 为行随机矩阵，\(P_{ij}\) 表示从状态
\(i\) 到状态 \(j\) 的一步转移概率。假设有严格正的平稳分布
\(\pi\)，作为列向量满足
\(\pi^\top P=\pi^\top\)、\(\mathbf1^\top\pi=1\)。记
\(D=\operatorname{diag}(\pi)\)。平稳时间反转为

\[
P^*=D^{-1}P^\top D.
\]

\end{minipage}\par

这一表达来自
\(\Pr(S_{t-1}=j\mid S_t=i)=\pi_jP_{ji}/\pi_i\)。只有平稳分布均匀时，转置才直接表示时间反转。无须假设过程非周期。

\par
\noindent\begin{minipage}{\linewidth}

包含当前状态的完整继任表征及纯未来部分分别为

\[
M_\gamma(P)=(I-\gamma P)^{-1}
=\sum_{h\ge0}\gamma^hP^h,
\qquad
N_\gamma(P)=P(I-\gamma P)^{-1},
\qquad 0\le\gamma<1.
\]

\end{minipage}\par

二者满足 \(M_\gamma=I+\gamma N_\gamma\)。因此 \(M_0=I\) 无方向信息，但
\(N_0=P\) 通常有一步方向信息。本文始终保留这一区别。

\par
\noindent\begin{minipage}{\linewidth}

令 \(\widetilde P=D^{1/2}PD^{-1/2}\)、\(r=\sqrt\pi\)，并取列标准正交的
\(U\in\mathbb R^{n\times(n-1)}\)，使 \(U^\top r=0\)。定义

\[
A=U^\top\widetilde P U=S+K,
\qquad S=(A+A^\top)/2,
\qquad K=(A-A^\top)/2.
\]

\end{minipage}\par

\(S\) 是对称部分，\(K\) 是相对于任务方向的反对称部分。令
\(F_\gamma(A)=(I-\gamma A)^{-1}\)，\(\operatorname{skew}T=(T-T^\top)/2\)，\(\langle V,W\rangle_F=\operatorname{tr}(V^\top W)\)。在二维旋转块中统一采用

\[
J=\begin{pmatrix}0&-1\\1&0\end{pmatrix},
\qquad J^2=-I,
\qquad J^\top=-J.
\]

\Needspace{6\baselineskip}
\subsubsection{命题
T01：平稳对比与共同测量关系}\label{ux547dux9898-t01ux5e73ux7a33ux5bf9ux6bd4ux4e0eux5171ux540cux6d4bux91cfux5173ux7cfb}

\par
\noindent\begin{minipage}{\linewidth}

有
\(\|\widetilde P\|_2\le1\)、\(\|A\|_2\le1\)、\(-I\preceq S\preceq I\)。若条件
\(c\) 的总体状态模式满足

\[
Y_c=g_c M_c(P)B+\mathbf1b_c^\top,
\qquad g_c>0,
\]

\end{minipage}\par

\par
\noindent\begin{minipage}{\linewidth}

其中 \(B\) 为各条件共用的任意线性特征映射，\(b_c\)
是状态无关偏置，则其加权状态对比满足

\[
X_c=U^\top D^{1/2}Y_c=g_c M_c(A)R,
\qquad R=U^\top D^{1/2}B.
\]

\end{minipage}\par

\par
\noindent\begin{minipage}{\linewidth}

\textbf{证明。} Jensen 不等式给出

\[
\sum_i\pi_i\Big(\sum_jP_{ij}x_j\Big)^2
\le\sum_{ij}\pi_iP_{ij}x_j^2
=\sum_j\pi_jx_j^2.
\]

\end{minipage}\par

故 \(\widetilde P\) 在欧氏范数下为压缩算子。又 \(\widetilde P r=r\) 与
\(\widetilde P^\top r=r\)，所以 \(r^\perp\)
是共同约化子空间；限制后的算子为 \(A\)，反转为
\(A^\top\)。有理函数限制到该子空间与先限制再取函数相同，状态无关偏置被
\(U^\top D^{1/2}\mathbf1=0\) 消去。对称部分的谱界来自
\(\|A\|_2\le1\)。证毕。

\textbf{适用边界。} 固定记录通道或成功解码不能独自证明生物学映射 \(B\)
在条件间相同。若允许任意 \(B_c\)，任意可逆候选 \(M_c\) 均可令
\(B_c=g_c^{-1}M_c^{-1}(Y_c-\mathbf1b_c^\top)\)
来解释观测。非线性分类概率也不自动满足上述线性关系。数学中的 \(P\)
是候选认知机制使用的转移；它等于任务生成器仍须明示或独立约束。

\Needspace{7\baselineskip}
\subsection{A2.
几何与参考造成的精确等价}\label{a2.-ux51e0ux4f55ux4e0eux53c2ux8003ux9020ux6210ux7684ux7cbeux786eux7b49ux4ef7}

\Needspace{6\baselineskip}
\subsubsection{定理
T02：规范几何的正规性边界}\label{ux5b9aux7406-t02ux89c4ux8303ux51e0ux4f55ux7684ux6b63ux89c4ux6027ux8fb9ux754c}

\par
\noindent\begin{minipage}{\linewidth}

若 \(P\) 双随机且 \(0<\gamma<1\)，则 \(M_\gamma(P)\) 与
\(M_\gamma(P^\top)\) 的完整平方欧氏行距离矩阵相同，当且仅当 \(P\)
正规，即 \(PP^\top=P^\top P\)。令
\(G_f=M_\gamma(P)M_\gamma(P)^\top\)、\(G_b=M_\gamma(P)^\top M_\gamma(P)\)，还有

\[
G_f-G_b=\gamma^2G_f(PP^\top-P^\top P)G_b,
\]

\end{minipage}\par

\[
\frac{\gamma^2}{(1+\gamma)^4}\|PP^\top-P^\top P\|_F
\le\|G_f-G_b\|_F
\le\frac{\gamma^2}{(1-\gamma)^4}\|PP^\top-P^\top P\|_F.
\]

\textbf{证明。} 两个 Gram 矩阵的逆之差为
\(\gamma^2(PP^\top-P^\top P)\)。用 \(X-Y=X(Y^{-1}-X^{-1})Y\)
得第一式。两者均将 \(\mathbf1\) 映为
\((1-\gamma)^{-2}\mathbf1\)，故其差已中心化。平方距离确定中心化 Gram
矩阵，因此完整距离相等当且仅当 \(G_f-G_b=0\)；可逆性和 \(\gamma>0\)
随即给出正规性。由于
\(\sigma_{min}(I-\gamma P)\ge1-\gamma\)、\(\|I-\gamma P\|_2\le1+\gamma\)，对上述乘积应用奇异值上下界即得范数界。证毕。

\par
\noindent\begin{minipage}{\linewidth}

\textbf{非均匀扩展。} 将行模式改为 \(M_\gamma(P)D^{-1/2}\) 与
\(M_\gamma(P^*)D^{-1/2}\)，即明确采用特征度量 \(D^{-1}\)，结论变为
\(\widetilde P\) 正规的充要条件。令

\[
\Delta=F_\gamma(\widetilde P)F_\gamma(\widetilde P)^\top-
F_\gamma(\widetilde P)^\top F_\gamma(\widetilde P),
\qquad G=D^{-1/2}\Delta D^{-1/2}.
\]

\end{minipage}\par

有 \(G\pi=0\)。若距离相同，对称矩阵 \(G\) 必为
\(\mathbf1v^\top+v\mathbf1^\top\)，从而
\(v+(v^\top\pi)\mathbf1=0\)。再左乘 \(\pi^\top\) 得
\(v=0\)。故距离相等等价于
\(\Delta=0\)，之后同前。未经加权的非均匀情况不由此证明。

\par
\noindent\begin{minipage}{\linewidth}

\textbf{纯未来扩展。} 对 \(N_\gamma(P)\)，同一规范几何结论在
\(\gamma=0\) 也成立。行列和相同保证距离相等等价于 \(N_\gamma\)
正规。正规 \(P\) 的有理函数正规；反向用

\[
P=N_\gamma(I+\gamma N_\gamma)^{-1}
\]

\end{minipage}\par

及谱定理证明，不要求 \(P\) 可逆。相反，完整 \(M_0=I\) 的等距不蕴含 \(P\)
正规。

这些结论针对规范特征几何、相同折扣和完整距离，不能直接推广到任意各向异性读出
\(M_\gamma B\)、相关距离、部分状态对或两侧不同折扣的拟合。

\Needspace{6\baselineskip}
\subsubsection{推论
T03：高度有向但几何相同的图族}\label{ux63a8ux8bba-t03ux9ad8ux5ea6ux6709ux5411ux4f46ux51e0ux4f55ux76f8ux540cux7684ux56feux65cf}

\par
\noindent\begin{minipage}{\linewidth}

对至少三个状态的循环置换 \(R\)，令

\[
P=pR+qR^\top+(1-p-q)I,
\qquad p,q>0,\quad p+q\le1.
\]

\end{minipage}\par

它是正规双随机矩阵；当 \(p\ne q\) 时不可逆，平稳熵产生为
\((p-q)\log(p/q)\)，在固定 \(p\)、\(q\downarrow0\)
时无界，但所有折扣下的前后规范几何相同。

\par
\noindent\begin{minipage}{\linewidth}

更一般地，在 \(\mathbf1^\perp\) 中取正交向量对 \((u_\ell,v_\ell)\)，令

\[
K_\sigma=\sum_{\ell=1}^k\sigma_\ell\epsilon_\ell
(u_\ell v_\ell^\top-v_\ell u_\ell^\top),
\quad P_\sigma=\mathcal J_n+s(I-\mathcal J_n)+K_\sigma,
\quad\mathcal J_n=\mathbf1\mathbf1^\top/n,
\]

\end{minipage}\par

其中 \(0\le s<1\)、\(\sigma_\ell\in\{-1,1\}\)、\(\epsilon_\ell>0\)
足够小以保证逐项非负。则 \(2^k\) 个已标记转移矩阵在所有折扣下等距。

\par
\noindent\begin{minipage}{\linewidth}

\textbf{证明。}
循环中所有项可交换，故正规；熵产生直接按相反有向边配对求和。第二个构造中各平面正交，故

\[
K_\sigma^2=-\sum_\ell\epsilon_\ell^2
(u_\ell u_\ell^\top+v_\ell v_\ell^\top)
\]

\end{minipage}\par

\par
\noindent\begin{minipage}{\linewidth}

与符号无关；在状态对比空间中的 Gram 矩阵为

\[
\big[(1-\gamma s)^2I-\gamma^2K_\sigma^2\big]^{-1}.
\]

\end{minipage}\par

常数方向也相同，结论成立。证毕。该计数不是未标记非同构图计数，也不保证随着维数增加仍有统一的正分离距离。

\Needspace{6\baselineskip}
\subsubsection{命题
T04：直接参考的方向正性与已学习参考的完整别名}\label{ux547dux9898-t04ux76f4ux63a5ux53c2ux8003ux7684ux65b9ux5411ux6b63ux6027ux4e0eux5df2ux5b66ux4e60ux53c2ux8003ux7684ux5b8cux6574ux522bux540d}

\par
\noindent\begin{minipage}{\linewidth}

若 \(X_0=g_0R\) 满行秩，\(X_1=g_1F_\gamma(A)R\)，则
\(T=X_1X_0^+=(g_1/g_0)F_\gamma(A)\)。独立转移参考 \(P_0=\mathbf1q^\top\)
的已学习完整 SR 在状态对比后也具有 \(g_0R\) 形式。令 \(K\ne0\)，定义

\[
\omega(T)=\frac{\langle T,K\rangle_F}{\|K\|_F^2}.
\]

\end{minipage}\par

\par
\noindent\begin{minipage}{\linewidth}

对直接完整 SR，\(\gamma>0\)
时前向为正、后向为负。对同一个图的两个已学习前向参考，

\[
T=aF_\gamma(A)(I-\eta A)
=a[I+(\gamma-\eta)AF_\gamma(A)],
\]

\end{minipage}\par

\par
\noindent\begin{minipage}{\linewidth}

有

\[
\omega(T)=a(\gamma-\eta)\rho_\gamma/\|K\|_F^2,
\qquad\rho_\gamma>0.
\]

\end{minipage}\par

\par
\noindent\begin{minipage}{\linewidth}

\textbf{证明。} 满行秩给出
\(X_0X_0^+=I\)。独立参考的非恒等部分各行相同，故在对比后消失。为证明正性，设

\[
H=I-\gamma S\succ0,\quad L=H^{-1/2}KH^{-1/2},\quad R_\gamma=(I-\gamma L)^{-1}.
\]

\end{minipage}\par

\par
\noindent\begin{minipage}{\linewidth}

则 \(F_\gamma(A)=H^{-1/2}R_\gamma H^{-1/2}\)，逆矩阵之差给出

\[
\operatorname{skew}F_\gamma=\gamma F_\gamma K F_\gamma^\top,
\qquad
\operatorname{skew}(AF_\gamma)=F_\gamma K F_\gamma^\top.
\]

\end{minipage}\par

\par
\noindent\begin{minipage}{\linewidth}

第二式在 \(\gamma=0\) 仍成立。迹的循环性及
\(R_\gamma^\top L R_\gamma=L(I-\gamma^2L^2)^{-1}\) 给出

\[
\rho_\gamma=\langle F_\gamma K F_\gamma^\top,K\rangle_F
=-\operatorname{tr}\{L^2(I-\gamma^2L^2)^{-1}\}>0.
\]

\end{minipage}\par

最后一步使用 \(-L^2\) 非零半正定，不要求 \(H\) 与 \(L\)
可交换。取内积即得所有符号公式。证毕。

\par
\noindent\begin{minipage}{\linewidth}

\textbf{完整对比算子的反例。} 七状态 Paley 转移取 \(P_{ij}=1/3\) 当
\(j-i\pmod7\in\{1,2,4\}\)。其六维状态对比块满足

\[
A=-I/6+K,\qquad K^2=-7I/36.
\]

\end{minipage}\par

\par
\noindent\begin{minipage}{\linewidth}

由 \(F_t(A)=\{(1+t/6)I+tK\}/\{(1+t/6)^2+7t^2/36\}\) 可直接验算

\[
F_{3/10}(A)F_{4/5}(A)^{-1}
=\frac{53}{48}I-\frac{25}{56}K
=\frac{634}{519}F_{75/173}(A^\top).
\]

\end{minipage}\par

因此两个真正前向但时域不同的已学习模式，其相对算子严格等于一个合法后向完整
SR
乘正增益。等价发生在全部六维对比算子上，超出了欧氏几何；常数状态方向明确不在该等价中。确定性均值相等也不自动意味着噪声分布相等。

\Needspace{7\baselineskip}
\subsection{A3.
连续族为何改变可识别的科学主张}\label{a3.-ux8fdeux7eedux65cfux4e3aux4f55ux6539ux53d8ux53efux8bc6ux522bux7684ux79d1ux5b66ux4e3bux5f20}

\Needspace{6\baselineskip}
\subsubsection{命题
T05：观测相容集合上的性质与方向锚定}\label{ux547dux9898-t05ux89c2ux6d4bux76f8ux5bb9ux96c6ux5408ux4e0aux7684ux6027ux8d28ux4e0eux65b9ux5411ux951aux5b9a}

令 \(\Phi:\Theta\to\mathcal Y\)
是包含所有合法机制和干扰参数的观测映射。性质 \(q(\theta)\)
能由无噪声观测确定，当且仅当它在每个非空纤维 \(\Phi^{-1}(y)\) 上为常量。

\textbf{证明。}
若两个参数映到同一观测却性质不同，任何只以该观测为输入的规则都不能同时正确；若各纤维上恒定，则定义
\(\widetilde q(y)=q(\theta)\)，其中任取
\(\theta\in\Phi^{-1}(y)\)，便得到良定义的观测性质。证毕。该命题是标准可辨识性原则；本研究的数学内容在于具体计算
SR 观测族与其纤维，而非重新命名这一原则。

\par
\noindent\begin{minipage}{\linewidth}

若允许的转移集合 \(\mathcal P\) 对平稳时间反转封闭，则

\[
\{M_\gamma(P):P\in\mathcal P\}
=\{M_\gamma(P^*):P\in\mathcal P\}.
\]

\end{minipage}\par

这是因为
\(P^{**}=P\)，两边只是重新编号。若需要方向主张，学得转移必须由设计或独立资料约束。若行为参数集合的漏盖概率至多
\(\beta\)，神经集合的漏盖概率至多 \(\alpha\)，则二者同时正确的概率至少
\(1-\alpha-\beta\)；这是事件并集界，不要求两类误差独立。

\Needspace{6\baselineskip}
\subsubsection{定理 T06：单平面直接广义 SR
的半平面塌缩}\label{ux5b9aux7406-t06ux5355ux5e73ux9762ux76f4ux63a5ux5e7fux4e49-sr-ux7684ux534aux5e73ux9762ux584cux7f29}

\par
\noindent\begin{minipage}{\linewidth}

取单位圆上的正交块

\[
A_\sigma=aI+\sigma\beta J,\quad a^2+\beta^2=1,\quad\beta>0,\quad\sigma\in\{-1,1\},
\]

\end{minipage}\par

\par
\noindent\begin{minipage}{\linewidth}

并规定合格直接参考。设

\[
G_\sigma=cI+\alpha N_\gamma(A_\sigma),
\quad c\in\mathbb R,\quad\alpha\ge0,\quad0\le\gamma<1.
\]

\end{minipage}\par

\par
\noindent\begin{minipage}{\linewidth}

则

\[
G_\sigma=uI+vJ,\quad
D_\gamma=1-2a\gamma+\gamma^2>0,
\]

\end{minipage}\par

\[
u=c+\alpha\frac{a-\gamma}{D_\gamma},
\qquad v=\sigma\alpha\frac{\beta}{D_\gamma}.
\]

该连续族的像恰为 \(\sigma v\ge0\) 的整个闭半平面；若严格要求
\(\alpha>0\)，像为开半平面 \(\sigma v>0\)。给定任意非零参考 \(b\)，观察
\(y=G_\sigma b\) 可唯一识别 \(u,v\) 和净旋转符号，但无法识别
\(\gamma\)。

\par
\noindent\begin{minipage}{\linewidth}

\textbf{证明。} 用
\((I-\gamma A_\sigma)^{-1}=\{(1-a\gamma)I+\sigma\beta\gamma J\}/D_\gamma\)
相乘得系数。对任意合法 \((u,v)\) 及任意 \(\gamma\)，令

\[
\alpha=|v|D_\gamma/\beta,
\qquad c=u-|v|(a-\gamma)/\beta,
\]

\end{minipage}\par

\par
\noindent\begin{minipage}{\linewidth}

即可得到同一算子，证明整个像及折扣别名。由于 \(b,Jb\) 正交等长，

\[
u=\frac{b^\top y}{\|b\|^2},\qquad
v=\frac{(Jb)^\top y}{\|b\|^2}
=\frac{\det[b,y]}{\|b\|^2}.
\]

\end{minipage}\par

这也说明两半平面并集对单个非零参考的任意目标均值均可零残差拟合。证毕。

\par
\noindent\begin{minipage}{\linewidth}

\textbf{扩展与边界。} 该单平面均值关系不提供检验单一几何 SR
核的剩余限制。若允许前后纯未来项非负混合，净系数为

\[
v=\beta\left(\frac{\alpha_F}{D_{\gamma_F}}-
\frac{\alpha_B}{D_{\gamma_B}}\right).
\]

\end{minipage}\par

两侧权重均可严格为正而保持任意规定 \(v\)，再用自由 \(c\) 匹配
\(u\)。所以正符号不证明组分纯前向。限制
\(c\ge0\)、增益范围或跨块共享参数时，族会缩小，必须重新求相容集合。

\Needspace{6\baselineskip}
\subsubsection{反例
T07：自由当前项使两学习条件完全换向而观测不变}\label{ux53cdux4f8b-t07ux81eaux7531ux5f53ux524dux9879ux4f7fux4e24ux5b66ux4e60ux6761ux4ef6ux5b8cux5168ux6362ux5411ux800cux89c2ux6d4bux4e0dux53d8}

\par
\noindent\begin{minipage}{\linewidth}

在任务块 \(J,-J\) 中，设 \(\gamma=1/2\)、\(\alpha=5\)。前向机制取当前项
\(c_F=1\)：

\[
G^F_+=-I+4J,\qquad G^F_-=-I-4J.
\]

\end{minipage}\par

\par
\noindent\begin{minipage}{\linewidth}

后向机制取 \(c_B=3\)：

\[
G^B_+=I-4J=-G^F_+,\qquad
G^B_-=I+4J=-G^F_-.
\]

\end{minipage}\par

\par
\noindent\begin{minipage}{\linewidth}

将未知共同基础映射从 \(B\) 改为 \(B'=-B\)，就有
\(G^F_\pm B=G^B_\pm B'\)。完整的两个条件模式及相对算子均相同：

\[
T_{\rm rel}=(-15I-8J)/17.
\]

\end{minipage}\par

\textbf{证明。}
\(N_{1/2}(J)=(-\tfrac12I+J)/(1+\tfrac14)\)，逐项代入并相乘即可。所有当前项和未来增益在例中均为正。证毕。

本反例允许共同映射未知且模型间不同；若独立合格参考已固定同一个
\(B\)，则不能采用 \(B'=-B\)，T06 的直接方向结果仍成立。它也不否定“完整
SR 加非负额外当前项”的定理：后者把当前与未来权重约束在不同的连续族中。

\Needspace{7\baselineskip}
\subsection{A4.
用互反任务解除参照歧义}\label{a4.-ux7528ux4e92ux53cdux4efbux52a1ux89e3ux9664ux53c2ux7167ux6b67ux4e49}

\Needspace{6\baselineskip}
\subsubsection{定理
T08：一般实压缩块的互反符号}\label{ux5b9aux7406-t08ux4e00ux822cux5b9eux538bux7f29ux5757ux7684ux4e92ux53cdux7b26ux53f7}

\par
\noindent\begin{minipage}{\linewidth}

设 \(A=S+K\) 为实压缩算子，\(K\ne0\)。两个任务在同一对比坐标中为
\(A,A^\top\)。若两条件均为完整前向
SR，允许不同折扣与正增益，其精确相对算子为

\[
T=aF_\gamma(A)(I-\eta A^\top),
\qquad a>0,\quad0\le\gamma,\eta<1,\quad\gamma+\eta>0.
\]

\end{minipage}\par

则 \(\langle T,K\rangle_F>0\)。两条件均后向时为负；以同一个对称部分
\(S\) 构造的对称机制为零。

\par
\noindent\begin{minipage}{\linewidth}

\textbf{证明。} 记 \(F=F_\gamma(A)\)，插入 \(F^{-1}\) 与 \(F^{-\top}\)
得

\[
\operatorname{skew}T/a
=F\{(\gamma+\eta)K-\gamma\eta(SK+KS)\}F^\top.
\]

\end{minipage}\par

\par
\noindent\begin{minipage}{\linewidth}

沿用 T04 中 \(H=I-\gamma S\)、\(L=H^{-1/2}KH^{-1/2}\)，令

\[
V=(I-\gamma^2L^2)^{-1},
\qquad Z=-L^2V\succeq0.
\]

\end{minipage}\par

\(K\ne0\) 保证 \(Z\ne0\)。因 \(S\) 与 \(H\) 可交换，花括号可写为
\(H^{1/2}L_0H^{1/2}\)，其中

\[
L_0=(\gamma+\eta)L-\gamma\eta(SL+LS).
\]

\par
\noindent\begin{minipage}{\linewidth}

迹的循环性给出

\[
\begin{aligned}
\langle T,K\rangle_F/a
&=-\operatorname{tr}(R_\gamma L_0R_\gamma^\top L)\\
&=-\operatorname{tr}(L_0LV)\\
&=\operatorname{tr}\{[(\gamma+\eta)I-2\gamma\eta S]Z\}.
\end{aligned}
\]

\end{minipage}\par

\par
\noindent\begin{minipage}{\linewidth}

最后一步使用 \(L\) 与 \(V\) 可交换；没有假设 \(S\) 与 \(L\) 可交换。又

\[
(\gamma+\eta)I-2\gamma\eta S
\succeq[\gamma(1-\eta)+\eta(1-\gamma)]I\succ0.
\]

\end{minipage}\par

正定矩阵与非零半正定矩阵乘积的迹严格为正。后向以 \(A^\top=S-K\) 替换
\(A\)，使标记反号；对称机制的两个因子为同一个对称矩阵的函数，故乘积对称。证毕。

\textbf{零边界。} 仅一个折扣为零仍有严格方向；两者均为零时
\(T=aI\)。\(K=0\)
时任务本身没有方向标记。允许折扣、非可逆性或增益趋零时，严格符号不等于统一正的噪声分离距离。该定理也未包括条件间切换方向的
FB、BF 机制；这些机制在 T13 明确进入竞争族。

\Needspace{6\baselineskip}
\subsubsection{命题
T09：非负额外当前成分的已证扩展}\label{ux547dux9898-t09ux975eux8d1fux989dux5916ux5f53ux524dux6210ux5206ux7684ux5df2ux8bc1ux6269ux5c55}

\par
\noindent\begin{minipage}{\linewidth}

设 \(r_0,s_0\ge0\)，考虑

\[
T=a[F_\gamma(A)+r_0I][F_\eta(A^\top)+s_0I]^{-1}.
\]

\end{minipage}\par

对二维实压缩块以及任意维正规块，T08 的方向符号保持；对满足 \(A^2=0\)
的任意实块也保持。

\par
\noindent\begin{minipage}{\linewidth}

\textbf{二维证明。} 正交对角化对称部分后，可写

\[
A=\begin{pmatrix}s_1&k\\-k&s_2\end{pmatrix},\quad
h_i=1-\gamma s_i>0,\quad j_i=1-\eta s_i>0.
\]

\end{minipage}\par

\par
\noindent\begin{minipage}{\linewidth}

令 \(d_\gamma=h_1h_2+\gamma^2k^2\)、\(d_\eta=j_1j_2+\eta^2k^2\)，

\[
\Delta=(j_2+s_0d_\eta)(j_1+s_0d_\eta)+\eta^2k^2>0.
\]

\end{minipage}\par

\par
\noindent\begin{minipage}{\linewidth}

显式求逆可得 \(\operatorname{skew}T=cK\)，其中

\[
c=\frac{a d_\eta}{2d_\gamma\Delta}
\{\eta(h_1+h_2+2r_0d_\gamma)+
\gamma(j_1+j_2+2s_0d_\eta)\}>0.
\]

\end{minipage}\par

\par
\noindent\begin{minipage}{\linewidth}

所有项非负且
\(\gamma+\eta>0\)，证明成立。正规实矩阵正交分解为一维实块与二维旋转块，逐块加和即可。若
\(A^2=0\)，直接展开得

\[
\operatorname{skew}T
=a\left\{\frac{\gamma}{1+s_0}
+\frac{(1+r_0)\eta}{(1+s_0)^2}\right\}K.
\]

\end{minipage}\par

其余项仅为 \(I\) 与 \(AA^\top\)
的倍数，因而对称。证毕。此处没有声称任意高维非正规算子、任意独立特征映射或自由当前/未来比均受同一证明覆盖。

\Needspace{7\baselineskip}
\subsection{A5.
满秩要求与任务相容测量}\label{a5.-ux6ee1ux79e9ux8981ux6c42ux4e0eux4efbux52a1ux76f8ux5bb9ux6d4bux91cf}

\Needspace{6\baselineskip}
\subsubsection{定理
T10：任意算子功能量的精确可观测条件}\label{ux5b9aux7406-t10ux4efbux610fux7b97ux5b50ux529fux80fdux91cfux7684ux7cbeux786eux53efux89c2ux6d4bux6761ux4ef6}

\par
\noindent\begin{minipage}{\linewidth}

设 \(X_1=TX_0\)，暂不限制 \(T\) 的参数形式，令
\(\Pi=X_0X_0^+\)。数据确定功能量 \(\langle T,W\rangle_F\)，当且仅当

\[
W=W\Pi.
\]

\end{minipage}\par

\par
\noindent\begin{minipage}{\linewidth}

\textbf{证明。} 数据精确确定 \(T\Pi=X_1X_0^+\)。若 \(W=W\Pi\)，则
\(\langle T\Pi,W\rangle_F=\langle T,W\rangle_F\)。反之，取
\(\Delta=W(I-\Pi)\ne0\)。有 \(\Delta X_0=0\)，但

\[
\langle\Delta,W\rangle_F=\|W(I-\Pi)\|_F^2>0.
\]

\end{minipage}\par

因此 \(T\) 与 \(T+\Delta\) 观测相同而功能量不同。证毕。

完整任意算子要求 \(X_0\) 满行秩；若有 \(n\)
个状态并去掉常数方向，线性通道数 \(p\ge n-1\)
只是满秩的必要计数条件，并不保证模式线性独立或稳定。恢复指定标记只需标记的相关方向可见。受限参数模型还可进一步减弱要求，如
T12 与 T17；不能把任意算子的必要条件当成所有机制性质的必要条件。

\Needspace{6\baselineskip}
\subsubsection{\texorpdfstring{命题 T11：共同约化块与八循环的 \(P/P^3\)
干预}{命题 T11：共同约化块与八循环的 P/P\^{}3 干预}}\label{ux547dux9898-t11ux5171ux540cux7ea6ux5316ux5757ux4e0eux516bux5faaux73afux7684-pp3-ux5e72ux9884}

\par
\noindent\begin{minipage}{\linewidth}

若 \(E^\top E=I\) 且

\[
AE=EA_E,\qquad A^\top E=EA_E^\top,
\]

\end{minipage}\par

则对在相关谱上合法的有理函数 \(f\) 有
\(E^\top f(A)=f(A_E)E^\top\)。所以可先取共同任务块再分析
SR，而不引入其余块的未知作用。

\textbf{证明。} 两个关系表示子空间及正交补均保持不变。\(A\)
在相应正交基中为块对角矩阵；求幂、求可逆矩阵的逆和取有理函数均保持相同分块，结论成立。证毕。任意状态投影不自动具有该性质。

\par
\noindent\begin{minipage}{\linewidth}

对于八状态行循环 \(P_{j,j+1\bmod8}=1\)、\(Q=P^3\)，定义

\[
C_{1j}=\tfrac12\cos(\pi j/2),\qquad
C_{2j}=-\tfrac12\sin(\pi j/2),\qquad j=0,\ldots,7.
\]

\end{minipage}\par

\par
\noindent\begin{minipage}{\linewidth}

直接相乘有

\[
C\mathbf1=0,\quad CC^\top=I,\quad
CP=JC,\quad CQ=-JC,\quad
CP^\top=-JC,\quad CQ^\top=JC.
\]

\end{minipage}\par

\par
\noindent\begin{minipage}{\linewidth}

因此整图虽非互为反转，\(Q\ne P^\top\)，这个预定平面上却严格互反。对于完整
SR，

\[
T_E=a(I-\gamma J)^{-1}(I+\eta J)
=\frac{a}{1+\gamma^2}
\{(1-\gamma\eta)I+(\gamma+\eta)J\}.
\]

\end{minipage}\par

\par
\noindent\begin{minipage}{\linewidth}

完整标记

\[
W=\tfrac12(P-P^\top-Q+Q^\top)=C^\top(2J)C
\]

\end{minipage}\par

\par
\noindent\begin{minipage}{\linewidth}

的秩为二、\(\|W\|_F^2=8\)，且

\[
\frac{\langle a(I-\gamma P)^{-1}(I-\eta Q),W\rangle_F}{8}
=\frac{a(\gamma+\eta)}{2(1+\gamma^2)}.
\]

\end{minipage}\par

\textbf{恒等式验证。} 以循环幂展开
\((I-\gamma P)^{-1}=(\sum_{j=0}^7\gamma^jP^j)/(1-\gamma^8)\)。各幂在
Frobenius 内积下正交，标记只提取第 1、7、3、5 次幂系数。合并分子得到
\((\gamma+\eta)(1-\gamma^2+\gamma^4-\gamma^6)\)，约去
\((1-\gamma^2)(1+\gamma^4)\) 即得公式。\(W=C^\top(2J)C\)
亦可逐项验证。后向转置使符号反转；两个可交换的对称化转移给零分数。证毕。

该命题是任务结构结论，不宣称任何已采集参与者的共同映射、可见性或噪声条件已经合格。

\Needspace{6\baselineskip}
\subsubsection{定理
T12：任务相容读出与弱可见性}\label{ux5b9aux7406-t12ux4efbux52a1ux76f8ux5bb9ux8bfbux51faux4e0eux5f31ux53efux89c1ux6027}

\par
\noindent\begin{minipage}{\linewidth}

对训练后冻结的非零模板 \(R\in\mathbb R^{2\times p}\)，定义

\[
z_R(X)=\frac1{\|R\|_F}
\begin{pmatrix}\langle R,X\rangle_F\\
\langle JR,X\rangle_F\end{pmatrix}.
\]

\end{minipage}\par

\par
\noindent\begin{minipage}{\linewidth}

则对所有 \(X,u,v\) 有

\[
z_R(JX)=Jz_R(X),\qquad
z_R((uI+vJ)X)=(uI+vJ)z_R(X).
\]

\end{minipage}\par

\par
\noindent\begin{minipage}{\linewidth}

若两个总体模式满足 \(X_+=T X_-\)、\(T=uI+vJ\)，则压缩均值
\(x_\pm=z_R(X_\pm)\) 满足相同关系，并且

\[
\det[x_-,x_+]=v\|x_-\|^2.
\]

\end{minipage}\par

所以只要真实参考投影 \(x_-\ne0\)，便可识别 \(v\)
的符号；不需要两行参考满秩。

\textbf{证明。} \(J^\top=-J\) 给出
\(\langle R,JX\rangle_F=-\langle JR,X\rangle_F\)、\(\langle JR,JX\rangle_F=\langle R,X\rangle_F\)，即交换关系。又
\(R,JR\) 正交等长，故 \(\|z_R(X)\|\le\|X\|_F\)。行列式公式来自
\(\det[x,Jx]=\|x\|^2\)。证毕。

这是一项结构恒等式，不要求训练模板等于真实参考；但模板非零不保证正式均值有非零投影。落在
\(\operatorname{span}\{R,JR\}\) 正交补的信号会丢失。任意线性压缩 \(L\)
若要保持候选作用，须另验
\(LT_\theta=\widetilde T_\theta L\)；方差最大化的状态 PCA
或分类器不自动满足这一条件。

\par
\noindent\begin{minipage}{\linewidth}

\textbf{一般旋转块的两学习条件接口。} 设
\(A=\rho I+\kappa J\)，\(\rho^2+\kappa^2\le1\)、\(\kappa>0\)。纯前向相对算子有

\[
u=\frac{a\{(1-\gamma\rho)(1-\eta\rho)-\gamma\eta\kappa^2\}}
{(1-\gamma\rho)^2+\gamma^2\kappa^2},
\]

\end{minipage}\par

\[
v=\frac{a\kappa(\gamma+\eta-2\gamma\eta\rho)}
{(1-\gamma\rho)^2+\gamma^2\kappa^2}>0.
\]

正性由 \(\gamma(1-\eta\rho)+\eta(1-\gamma\rho)>0\) 得到。一般块的 \(u\)
可以为负；只有进一步限制为下节 quarter-turn 块，才有特定的正 \(u\)
扇区。该接口允许参考本身已学习，但不能跨越 T07 中自由当前项造成的等价。

\Needspace{7\baselineskip}
\subsection{A6.
四种机制族、混合与部分定量识别}\label{a6.-ux56dbux79cdux673aux5236ux65cfux6df7ux5408ux4e0eux90e8ux5206ux5b9aux91cfux8bc6ux522b}

\Needspace{6\baselineskip}
\subsubsection{定理 T13：quarter-turn
相对连续族的精确角度像}\label{ux5b9aux7406-t13quarter-turn-ux76f8ux5bf9ux8fdeux7eedux65cfux7684ux7cbeux786eux89d2ux5ea6ux50cf}

\par
\noindent\begin{minipage}{\linewidth}

在任务块 \(J,-J\) 中，两个条件均使用纯完整 SR，各自折扣在
\([0,g]\)，其中 \(g<1\)，两条件正增益比任意。以 F、B
分别表示相对于各自任务的前、后向机制。令 \(h=\arctan g\)，将相对算子写为

\[
T=r(\cos\phi\,I+\sin\phi\,J),\qquad r>0.
\]

\end{minipage}\par

则四族的准确角度像为

\Needspace{9\baselineskip}

\begingroup
\small

\begin{longtable}[]{@{}ll@{}}
\toprule\noalign{}
机制标签 & 相对角 \(\phi\) 的全部范围 \\
\midrule\noalign{}
\endhead
\bottomrule\noalign{}
\endlastfoot
FF & \([0,2h]\) \\
BB & \([-2h,0]\) \\
FB & \([-h,h]\) \\
BF & \([-h,h]\) \\
\end{longtable}

\endgroup

每个范围中的每个角度和任意正径向尺度均可实现。

\par
\noindent\begin{minipage}{\linewidth}

\textbf{证明。} 单个完整 SR 满足

\[
F_\gamma(J)=\frac{I+\gamma J}{1+\gamma^2}
=\frac1{\sqrt{1+\gamma^2}}
\{\cos\theta_\gamma I+\sin\theta_\gamma J\},
\qquad\theta_\gamma=\arctan\gamma.
\]

\end{minipage}\par

\par
\noindent\begin{minipage}{\linewidth}

反任务上的前向角为
\(-\theta_\eta\)；相对算子的角是目标减参考，故四族分别为

\[
\theta_\gamma+\theta_\eta,\quad
-\theta_\gamma-\theta_\eta,\quad
\theta_\gamma-\theta_\eta,\quad
-\theta_\gamma+\theta_\eta.
\]

\end{minipage}\par

独立角度取值于 \([0,h]\)，即得区间。FF 中取
\(\theta_\gamma=\theta_\eta=\phi/2\)
实现每个角；切换族中取一个角为零、另一个取所需绝对值。增益比自由可匹配任意
\(r>0\)。证毕。

\par
\noindent\begin{minipage}{\linewidth}

FF 的角与尺度具体为

\[
\phi=\arctan\gamma+\arctan\eta,\qquad
r=a\sqrt{\frac{1+\eta^2}{1+\gamma^2}}.
\]

\end{minipage}\par

因此 \(0<\phi\le h\) 可由 FF 与切换机制同时产生，正面积不能单独识别
FF；若数据相容集合非空且全部满足 \(\phi>h\)，则只有 FF
相容。其可观测形式为 \(u>0\) 且
\(v-gu>0\)。零折扣和净方向零边界必须保留。若增益有额外上下界，径向范围不能继续当作全正轴。

当折扣范围为 \([0,1)\) 时，FF、BB 角区间分别为
\([0,\pi/2)\)、\((-\pi/2,0]\)，切换族为
\((-\pi/4,\pi/4)\)。因此精确无噪声的 \(\phi=\pi/4\)
已不属于切换族，但它属于该族闭包；带稳健余量的排除仍要求置信范围严格超过
\(\pi/4\)。

\Needspace{6\baselineskip}
\subsubsection{命题 T14：非负完整 SR
混合保留净方向扇区}\label{ux547dux9898-t14ux975eux8d1fux5b8cux6574-sr-ux6df7ux5408ux4fddux7559ux51c0ux65b9ux5411ux6247ux533a}

\par
\noindent\begin{minipage}{\linewidth}

在各自任务坐标中，设一个条件允许

\[
H(J)=cI+\int_0^g F_\gamma(J)\,\mu_F(d\gamma)
+\int_0^g F_\gamma(-J)\,\mu_B(d\gamma),
\]

\end{minipage}\par

\par
\noindent\begin{minipage}{\linewidth}

其中 \(c\ge0\)，\(\mu_F,\mu_B\)
为有限非负测度，且该算子非零；所有尺度、方向共用同一个特征映射。则

\[
H(J)=uI+vJ,\quad
u=c+\int_0^g\frac{\mu_F+\mu_B}{1+\gamma^2}>0,
\]

\end{minipage}\par

\[
v=\int_0^g\frac{\gamma(\mu_F-\mu_B)}{1+\gamma^2},
\qquad |v|\le gu.
\]

故每个条件净相位 \(\theta=\arctan(v/u)\) 均在
\([-h,h]\)，按两个条件净相位的正负分类仍得到 T13
的四个扇区。若相对角大于 \(h\)，两条件相对于各自任务都必须净前向。

\textbf{证明。} 展开每个完整 SR 得系数；逐项用 \(\gamma\le g\)
和非负权重得到不等式。两条件相对角是各自任务坐标净相位之和，区间运算同
T13。证毕。

此结论识别的是加权净方向，而非没有任何相反方向组分。它不涵盖任意自由
\(cI+\alpha N_\gamma\) 的相对族；后者可有 \(u<0\)，T07 已给反例。

\Needspace{6\baselineskip}
\subsubsection{定理
T15：从方向到折扣的尖锐部分识别}\label{ux5b9aux7406-t15ux4eceux65b9ux5411ux5230ux6298ux6263ux7684ux5c16ux9510ux90e8ux5206ux8bc6ux522b}

\par
\noindent\begin{minipage}{\linewidth}

在纯完整 SR 的 FF 模型中，设
\(\gamma,\eta\in[0,g]\)、\(h=\arctan g\)，给定精确相对角
\(\phi\in[0,2h]\)，且增益比自由。则

\[
\gamma\in
\left[
\tan\{\max(0,\phi-h)\},
\tan\{\min(h,\phi)\}
\right],
\]

\end{minipage}\par

\(\eta\) 具有相同边缘范围，联合相容集为
\(\arctan\gamma+\arctan\eta=\phi\)。每个允许值均可实现，因而这些界是尖锐的。

\par
\noindent\begin{minipage}{\linewidth}

\textbf{证明。} 令
\(\theta_1=\arctan\gamma\)、\(\theta_2=\arctan\eta\)。约束为
\(0\le\theta_i\le h\)、\(\theta_1+\theta_2=\phi\)，等价于

\[
\max(0,\phi-h)\le\theta_1\le\min(h,\phi).
\]

\end{minipage}\par

\(\tan\) 严格递增，得到边缘区间。对区间内任取 \(\theta_1\)，令
\(\theta_2=\phi-\theta_1\)，再按 T13
选择增益匹配观测尺度，即达到该值。证毕。

\par
\noindent\begin{minipage}{\linewidth}

若原范围为 \([0,1)\)，则有三个精确情形：

\[
\begin{array}{ll}
0\le\phi<\pi/4:
&\gamma,\eta\text{ 的边缘识别集为 }[0,\tan\phi],\\
\phi=\pi/4:
&\gamma,\eta\text{ 的边缘识别集为 }(0,1),\\
\pi/4<\phi<\pi/2:
&\gamma,\eta\text{ 的边缘识别集为 }
(\tan(\phi-\pi/4),1).
\end{array}
\]

\end{minipage}\par

\par
\noindent\begin{minipage}{\linewidth}

特别是

\[
\boxed{\phi>\pi/4\quad\Longrightarrow\quad
\gamma>\tan(\phi-\pi/4),\quad
\eta>\tan(\phi-\pi/4).}
\]

\end{minipage}\par

例如 \(\phi=70^\circ\) 强制两个折扣均大于
\(\tan25^\circ\approx0.4663\)，但通常不唯一决定任一个折扣。有限上界
\(g\) 下，\(\phi=0\) 或 \(2h\)
等边界可使两折扣唯一；“一般不可单独恢复”不能替代这些边界事实。

对 T14 的混合模型，同样的相位逻辑只约束有效净系数
\(q_i=v_i/u_i\)，不能逐个约束混合中的每个
\(\gamma\)。有限噪声下应投影整个合法联合相容集合；不能把点估计角代入上述公式后直接称作置信界。

\Needspace{6\baselineskip}
\subsubsection{命题
T16：单旋转平面不区分同方向多尺度与单一折扣}\label{ux547dux9898-t16ux5355ux65cbux8f6cux5e73ux9762ux4e0dux533aux5206ux540cux65b9ux5411ux591aux5c3aux5ea6ux4e0eux5355ux4e00ux6298ux6263}

\par
\noindent\begin{minipage}{\linewidth}

设
\(H=cI+\int F_\gamma(J)\mu(d\gamma)\)，\(c\ge0\)、\(\mu\ge0\)、\(H\ne0\)。写
\(H=uI+vJ\)。则

\[
\gamma_{\rm eff}=v/u,\qquad
g_{\rm eff}=u(1+\gamma_{\rm eff}^2)
\]

\end{minipage}\par

满足 \(H=g_{\rm eff}F_{\gamma_{\rm eff}}(J)\)。若所有折扣小于一，则
\(0\le\gamma_{\rm eff}<1\)。

\textbf{证明。} \(u>0\)、\(v\ge0\)，且 \(v/u\) 为以 \(1/(1+\gamma^2)\)
加权的折扣平均，并包含当前项的零折扣质量，故在合法范围内。直接代入
\(F_{\gamma_{\rm eff}}(J)=(I+\gamma_{\rm eff}J)/(1+\gamma_{\rm eff}^2)\)
即得。证毕。

更多同平面模板可以检验共享线性算子，但不能拆开已在该平面完全相同的时间核。定量出口需要额外传播形状或额外任务块。

\Needspace{7\baselineskip}
\subsection{A7.
识别折扣与时间核的建设性出口}\label{a7.-ux8bc6ux522bux6298ux6263ux4e0eux65f6ux95f4ux6838ux7684ux5efaux8bbeux6027ux51faux53e3}

\Needspace{6\baselineskip}
\subsubsection{定理
T17：三个传播形状是自由当前项模型的精确折扣条件}\label{ux5b9aux7406-t17ux4e09ux4e2aux4f20ux64adux5f62ux72b6ux662fux81eaux7531ux5f53ux524dux9879ux6a21ux578bux7684ux7cbeux786eux6298ux6263ux6761ux4ef6}

\par
\noindent\begin{minipage}{\linewidth}

设 \(P\in\mathbb R^{m\times m}\) 已知且谱半径不超过一，直接参考
\(B\in\mathbb R^{m\times p}\) 已知。考虑完整参数域

\[
Y(c,\alpha,\gamma)
=cB+\alpha P(I-\gamma P)^{-1}B,
\qquad c\in\mathbb R,\quad\alpha>0,\quad0\le\gamma<1.
\]

\end{minipage}\par

则以下命题等价：

\begin{enumerate}
\def\labelenumi{\arabic{enumi}.}
\tightlist
\item
  \(\gamma\) 在该族内全局结构可辨识。
\item
  \(c,\alpha,\gamma\) 三个参数在该族内共同可辨识。
\item
  \(B,PB,P^2B\) 作为同维矩阵在实标量线性组合意义下线性独立。
\end{enumerate}

若第三条不成立，任意给定合法参数点的相同观测都可由任意另一个合法折扣产生；不只是存在一对例外碰撞。

\par
\noindent\begin{minipage}{\linewidth}

\textbf{充分性证明。} 假设两组参数产生相同 \(Y\)，记 \(\delta\)
为第二组折扣，\(c',\alpha'\) 为其余参数。所有矩阵均为 \(P\)
的函数，故左乘 \((I-\gamma P)(I-\delta P)\) 得

\[
\begin{aligned}
0={}&(c-c')B\\
&+\{\alpha-\alpha'-(c-c')(\gamma+\delta)\}PB\\
&+\{(c-c')\gamma\delta-\alpha\delta+\alpha'\gamma\}P^2B.
\end{aligned}
\]

\end{minipage}\par

三个传播形状线性独立，系数分别为零。第一式给 \(c=c'\)，第二式给
\(\alpha=\alpha'\)，第三式变为 \(\alpha(\gamma-\delta)=0\)。因
\(\alpha>0\)，得 \(\gamma=\delta\)，证明共同唯一性。

\par
\noindent\begin{minipage}{\linewidth}

\textbf{必要性证明。} 若 \(B=0\)，所有观测为零。若 \(B\ne0\) 而
\(PB=\lambda B\)，则

\[
Y=\left(c+\frac{\alpha\lambda}{1-\gamma\lambda}\right)B.
\]

\end{minipage}\par

任意改变 \(\gamma\) 后，都可保持正 \(\alpha\) 并改变自由 \(c\)
使观测不变。

\par
\noindent\begin{minipage}{\linewidth}

余下情况为 \(B,PB\) 线性独立，但存在实数 \(s,t\) 使

\[
P^2B=sB+tPB.
\]

\end{minipage}\par

\par
\noindent\begin{minipage}{\linewidth}

于是 \(\operatorname{span}\{B,PB\}\) 对左乘 \(P\) 不变。在有序基
\((B,PB)\) 上，左乘 \(P\) 的矩阵为

\[
C=\begin{pmatrix}0&s\\1&t\end{pmatrix}.
\]

\end{minipage}\par

\par
\noindent\begin{minipage}{\linewidth}

其谱是左乘 \(P\) 在该子空间的谱；故 \(I-\gamma C\) 对 \(0\le\gamma<1\)
可逆。令

\[
D_\gamma=\det(I-\gamma C)=1-t\gamma-s\gamma^2.
\]

\end{minipage}\par

\(D_\gamma\) 在整个连通区间中连续、不为零，且
\(D_0=1\)，所以始终严格为正。直接求逆得

\[
P(I-\gamma P)^{-1}B
=\frac{\gamma s B+PB}{D_\gamma},
\]

\par
\noindent\begin{minipage}{\linewidth}

从而

\[
Y=\left(c+\frac{\alpha\gamma s}{D_\gamma}\right)B
+\frac{\alpha}{D_\gamma}PB.
\]

\end{minipage}\par

\par
\noindent\begin{minipage}{\linewidth}

对任意另一个折扣 \(\delta\in[0,1)\)，取

\[
\alpha'=\alpha D_\delta/D_\gamma>0,
\qquad
c'=c+\alpha s(\gamma-\delta)/D_\gamma.
\]

\end{minipage}\par

代入即有
\(Y(c',\alpha',\delta)=Y(c,\alpha,\gamma)\)，证明必要性和逐点别名。证毕。

\par
\noindent\begin{minipage}{\linewidth}

\textbf{解释与范围。}
这里的线性独立是将三个矩阵向量化后检查三列的秩，不是要求 \(B\)
行满秩。比如 \(J_3\) 为三阶下移位块、\(B=e_1\)，则

\[
B=e_1,\quad J_3B=e_2,\quad J_3^2B=e_3,\qquad
Y=(c,\alpha,\alpha\gamma)^\top.
\]

\end{minipage}\par

一列通道即可辨认三个参数，因为任务让三个传播形状分开。T19
将这一块嵌入合法随机任务；本例仅解释代数条件，不代表神经实验已成功。

本定理假设 \(P\) 和直接 \(B\) 已知、\(\alpha\) 严格为正、\(c\)
全实且无额外参数限制；改变这些条件需重新判断。若 \(\alpha=0\)
也允许，则该边界所有 \(\gamma\) 均同像。若 \(P\)
或参考未知，不能借本定理绕过 T05 或
T07。它是在指定方向分支内的参数识别，不独自识别机制方向。线性独立也不是有限噪声下稳定估计的保证；近相关的传播形状会使置信区间宽。

\Needspace{6\baselineskip}
\subsubsection{定理
T18：两个不同旋转块的显式折扣恢复}\label{ux5b9aux7406-t18ux4e24ux4e2aux4e0dux540cux65cbux8f6cux5757ux7684ux663eux5f0fux6298ux6263ux6062ux590d}

\par
\noindent\begin{minipage}{\linewidth}

采用 T17 的直接参考模型，设可见的两个任务块为

\[
A_j=a_jI+\sigma\beta_jJ,\quad
a_j^2+\beta_j^2=1,\quad\beta_j>0,\quad a_1\ne a_2,
\]

\end{minipage}\par

\par
\noindent\begin{minipage}{\linewidth}

且两个块共享同一
\(c,\alpha,\gamma\)，\(\alpha>0\)。每个参考均有非零相容投影，所以可观测
\(u_j,v_j\)。有

\[
D_j=1-2a_j\gamma+\gamma^2,\quad
u_j=c+\alpha(a_j-\gamma)/D_j,\quad
v_j=\sigma\alpha\beta_j/D_j.
\]

\end{minipage}\par

\par
\noindent\begin{minipage}{\linewidth}

令

\[
q=\frac{v_1/\beta_1}{v_2/\beta_2}=\frac{D_2}{D_1},
\qquad
t=\frac{q-1}{2(qa_1-a_2)}.
\]

\end{minipage}\par

\par
\noindent\begin{minipage}{\linewidth}

则合法模型中 \(t=\gamma/(1+\gamma^2)\) 且

\[
\boxed{\gamma=\frac{2t}{1+\sqrt{1-4t^2}}.}
\]

\end{minipage}\par

\par
\noindent\begin{minipage}{\linewidth}

随后

\[
\alpha=\frac{\sigma v_jD_j}{\beta_j},\qquad
c=u_j-\alpha(a_j-\gamma)/D_j.
\]

\end{minipage}\par

\par
\noindent\begin{minipage}{\linewidth}

\textbf{证明。} \(q=(1-2a_2t)/(1-2a_1t)\)，直接求解得 \(t\)。又

\[
\frac{dq}{d\gamma}
=\frac{2(a_1-a_2)(1-\gamma^2)}{D_1^2}
\]

\end{minipage}\par

在 \([0,1)\) 上符号严格固定，所以一一对应。\(t\) 严格递增于
\([0,1/2)\)，解其二次方程并选小于一的根，得所述逆式；写成有理化形式避免
\(\gamma=0\) 时的 \(0/0\)。再用两块系数恢复其余参数。证毕。

\textbf{剩余可检验约束。} 两块各自算出的 \(c\)
必须一致；这使四个可观测系数承担三个共享参数，增加一个均值层面的可失败限制。例如八循环的两个实际
Fourier 块可取
\((a_1,\beta_1)=(0,1)\)、\((a_2,\beta_2)=(\sqrt2/2,\sqrt2/2)\)。若
\(v_1=4/5\)、\(v_2=(8+10\sqrt2)/17\)，则强制
\(\gamma=1/2,\alpha=1\)；若又观察 \(u_1=u_2=1\)，两块给
\(c_1=7/5,c_2=(19-6\sqrt2)/17\)，模型不相容。这里明确选择有限随机任务中可实现的单位圆块；任意抽象正交块并不自动对应有限马尔可夫链的单位模特征值。

\par
\noindent\begin{minipage}{\linewidth}

也可令 \(s_j=\sigma v_j/\beta_j\)，在 \(s_1\ne s_2\) 时有

\[
\gamma=\frac{s_1a_1-s_2a_2-(u_1-u_2)}{s_1-s_2}.
\]

\end{minipage}\par

该式在 \(\gamma=0\) 时退化，而比值逆式仍可用。比值对 \(\gamma\) 的导数在
\(\gamma\to1\)
时趋零，只说明该比值的敏感性变差，不证明完整系数信息必同样退化。两个块若可独立选择增益、当前项或折扣，以上共同参数结论不再成立。统计估计应联合反演系数和参考误差，而非对不稳定比值直接套正态近似。

\Needspace{6\baselineskip}
\subsubsection{命题
T19：合法随机任务中的可扩展幂零块}\label{ux547dux9898-t19ux5408ux6cd5ux968fux673aux4efbux52a1ux4e2dux7684ux53efux6269ux5c55ux5e42ux96f6ux5757}

\par
\noindent\begin{minipage}{\linewidth}

以二元字 \(x=(x_1,\ldots,x_k)\)
为状态，每次删除第一位并在末尾加入独立等概率二元位。完整图有 \(n=2^k\)
个状态，每状态两条概率为 \(1/2\) 的入边与出边；故为双随机矩阵。定义

\[
C_{j,x}=n^{-1/2}(1-2x_j),\qquad j=1,\ldots,k.
\]

\end{minipage}\par

\par
\noindent\begin{minipage}{\linewidth}

有

\[
C\mathbf1=0,\quad CC^\top=I_k,\quad
CP=J_kC,\quad CP^\top=J_k^\top C,\quad J_k^k=0,
\]

\end{minipage}\par

其中 \((J_k)_{i+1,i}=1\)。

\textbf{证明。}
各二元位在全部状态中独立平衡，得到零均值和正交性。固定终点后，前驱第一位是独立等概率变量，其余位由终点前
\(k-1\) 位确定；故 \(CP\) 第一行零，第 \(j>1\) 行等于 \(C\) 的第 \(j-1\)
行。对后继求平均得到转置关系。移位 \(k\) 次为零。证毕。

幂零的是预定对比块，不是完整转移矩阵；完整任务仍为正常随机转移系统。该构造使用经典
de Bruijn 图结构，本文用其获得有限滞后上可分开的传播方向。

\Needspace{6\baselineskip}
\subsubsection{定理
T20：直接参考下的有限时间核与混合检验}\label{ux5b9aux7406-t20ux76f4ux63a5ux53c2ux8003ux4e0bux7684ux6709ux9650ux65f6ux95f4ux6838ux4e0eux6df7ux5408ux68c0ux9a8c}

\par
\noindent\begin{minipage}{\linewidth}

在 T19 块中，若合格直接参考支持所需算子量的恢复，目标时间核为
\(f(z)=\sum_h w_hz^h\)，则

\[
T=a\sum_{h=0}^{k-1}w_hJ_k^h+cI.
\]

\end{minipage}\par

\par
\noindent\begin{minipage}{\linewidth}

定义

\[
\psi_h=\frac1{k-h}\sum_{i=1}^{k-h}
(T_{i+h,i}-T_{i,i+h}),\qquad 1\le h<k.
\]

\end{minipage}\par

则 \(\psi_h=aw_h\)；后向核给 \(-aw_h\)；对称算子给零。完整单一 SR 的
\(w_h=\gamma^h\)，故在有非零分母时 \(\psi_{h+1}/\psi_h=\gamma\)，且
\(k\ge4\) 时有 \(\psi_2^2=\psi_1\psi_3\)。

\par
\noindent\begin{minipage}{\linewidth}

对具有共同映射的非负完整 SR 时域混合
\(w_h=\int_0^1\gamma^h\mu(d\gamma)\)，有

\[
w_2^2\le w_1w_3.
\]

\end{minipage}\par

\par
\noindent\begin{minipage}{\linewidth}

当 \(w_1>0\)，令 \(\nu(d\gamma)=\gamma\mu(d\gamma)/w_1\)，则

\[
\frac{w_1w_3-w_2^2}{w_1^2}
=\operatorname{Var}_{\nu}(\gamma).
\]

\end{minipage}\par

\textbf{证明。} \(J_k^h\) 仅在第 \(h\) 条下副对角线上为一，逐项得到
\(\psi_h\)。将 Cauchy–Schwarz 应用于 \(\gamma^{1/2},\gamma^{3/2}\)
得不等式；归一化后
\(\mathbb E_\nu\gamma=w_2/w_1\)、\(\mathbb E_\nu\gamma^2=w_3/w_1\)，给方差恒等式。证毕。

方差也可用 \(\psi_h\)
表示，整体正增益消去。零折扣质量对正滞后不可见，等式不能排除它。有限可见滞后既不能验证无限尾部，也不能唯一确定学习算法。若各尺度具有任意不同的特征映射，混合矩不等式不再自动适用。

\Needspace{6\baselineskip}
\subsubsection{定理
T21：未知反转参考下的时间核与双折扣恢复}\label{ux5b9aux7406-t21ux672aux77e5ux53cdux8f6cux53c2ux8003ux4e0bux7684ux65f6ux95f4ux6838ux4e0eux53ccux6298ux6263ux6062ux590d}

\par
\noindent\begin{minipage}{\linewidth}

在同一幂零块内，令目标和反向参考核分别为

\[
f(J_k)=\sum_{h=0}^{k-1}f_hJ_k^h,\qquad
q(J_k^\top)=\sum_{h=0}^{k-1}q_h(J_k^\top)^h,
\qquad q_0\ne0.
\]

\end{minipage}\par

\par
\noindent\begin{minipage}{\linewidth}

若所需相对算子量可观测，则

\[
T=a f(J_k)q(J_k^\top)^{-1}
\quad\Longrightarrow\quad
T_{h+1,1}=\frac{a}{q_0}f_h.
\]

\end{minipage}\par

\textbf{证明。} \(J_k^\top e_1=0\)，所以
\(q(J_k^\top)e_1=q_0e_1\)，逆矩阵将 \(e_1\) 映为 \(e_1/q_0\)。再乘
\(f(J_k)\) 得第一列公式。证毕。

这在一个共同尺度之外恢复目标的可见核，不要求参考为几何核；方向解释另需
\(a/q_0>0\)。当前项只改变 \(f_0\) 或
\(q_0\)，所以正滞后比例保持。\(q_0=0\) 时参考块奇异，不支持这个论证。

\par
\noindent\begin{minipage}{\linewidth}

纯完整 SR 给出

\[
T=a(I-\gamma J_k)^{-1}(I-\eta J_k^\top),
\qquad a=T_{11},\quad
\gamma=T_{21}/T_{11},\quad
\eta=-T_{12}/T_{11}.
\]

\end{minipage}\par

这些式子在 \(k\ge2\) 时成立。\(k\ge4\) 时第一列还须满足
\(T_{31}^2=T_{21}T_{41}\)。参数识别以纯 SR
族为条件，其他矩阵元素和竞争核仍须比较。

不能将 T20
的直接副对角线差量误用于这个相对算子：上下副对角线均受参考核影响。精确单一
SR 在 \(\gamma=0.6,\eta=0.8,a=1,k=4\) 下，误用会得到
\(\psi_2^2-\psi_1\psi_3=-0.18607104\)，错误地似乎支持多尺度混合。相对第一列是适用于该参考问题的量。其可观测性也应按
T10 或受限模型另证，不能自动由二维方向读出推得。

\Needspace{6\baselineskip}
\subsubsection{命题
T22：零方向不是唯一对称机制}\label{ux547dux9898-t22ux96f6ux65b9ux5411ux4e0dux662fux552fux4e00ux5bf9ux79f0ux673aux5236}

\par
\noindent\begin{minipage}{\linewidth}

在 \(J_4\) 块上，混合两方向的完整地图

\[
\tfrac12\{F_\gamma(J_4)+F_\gamma(J_4^\top)\}
\]

\end{minipage}\par

\par
\noindent\begin{minipage}{\linewidth}

具有常数对角线。而先对称化转移再取 SR 的

\[
F_\eta((J_4+J_4^\top)/2)
\]

\end{minipage}\par

\par
\noindent\begin{minipage}{\linewidth}

在 \(0<\eta<1\) 时满足

\[
[F_\eta]_{22}-[F_\eta]_{11}
=\frac{x^2}{1-3x^2+x^4}>0,\qquad x=\eta/2.
\]

\end{minipage}\par

\textbf{证明。}
幂零完整地图的所有对角线均为一，前一个结论立即成立。后一矩阵的三对角行列式为
\(1-3x^2+x^4\)，两个相关余子式为 \(1-x^2\) 与
\(1-2x^2\)，相减即得公式。证毕。正增益保留差异，额外标量当前项消去；只有零反对称分数无法辨别这两种机制。

\Needspace{7\baselineskip}
\subsection{A8.
时间历史与误差预算的边界}\label{a8.-ux65f6ux95f4ux5386ux53f2ux4e0eux8befux5deeux9884ux7b97ux7684ux8fb9ux754c}

\Needspace{6\baselineskip}
\subsubsection{定理
T23：记录历史的后瞻别名与独立探测控制}\label{ux5b9aux7406-t23ux8bb0ux5f55ux5386ux53f2ux7684ux540eux77bbux522bux540dux4e0eux72ecux7acbux63a2ux6d4bux63a7ux5236}

\par
\noindent\begin{minipage}{\linewidth}

若真实瞬时响应只有当前状态码 \(b(S_t)\)，但测量为

\[
y_t=\sum_{h\ge0}\lambda^h b(S_{t-h}),\qquad0\le\lambda<1,
\]

\end{minipage}\par

\par
\noindent\begin{minipage}{\linewidth}

在平稳马尔可夫链下有

\[
\mathbb E[y_t\mid S_t=i]=[M_\lambda(P^*)B]_i.
\]

\end{minipage}\par

\textbf{证明。} 条件过去分布为
\(\Pr(S_{t-h}=j\mid S_t=i)=[(P^*)^h]_{ij}\)，逐项取条件期望并求和。证毕。因而图反转也会反转该测量伪影；任务互反本身不排除历史滤波机制。

\par
\noindent\begin{minipage}{\linewidth}

若学习后的探测状态 \(Z_t\) 独立同分布，并独立于学习序列，已学习响应
\(r_c(z)\) 在测量期稳定，且记录滤波器绝对可和，则

\[
y_t=\sum_{h\ge0}h_h r_c(Z_{t-h})
\quad\Longrightarrow\quad
\mathbb E[y_t\mid Z_t=z]
=h_0r_c(z)+\sum_{h\ge1}h_h\mathbb E r_c(Z).
\]

\end{minipage}\par

后项与当前状态无关，故状态对比将其消去。这是条件均值层面的控制。若需要方向符号，当前有效增益还须与观测模型规定一致；零系数会毁掉可见性。不放回排列、禁止相邻重复、状态依赖适应、继续学习或不稳定测量均不满足这一简单证明。

\Needspace{6\baselineskip}
\subsubsection{定理
T24：完整算子恢复的确定性误差界}\label{ux5b9aux7406-t24ux5b8cux6574ux7b97ux5b50ux6062ux590dux7684ux786eux5b9aux6027ux8befux5deeux754c}

\par
\noindent\begin{minipage}{\linewidth}

设

\[
\widehat X_0=X_0+E_0,\quad
\widehat X_1=TX_0+E_1,\quad
\|E_i\|_F\le e_i.
\]

\end{minipage}\par

\par
\noindent\begin{minipage}{\linewidth}

若 \(\widehat X_0\) 满行秩且 \(s=\sigma_{\min}(\widehat X_0)>e_0\)，定义
\(\widehat T=\widehat X_1\widehat X_0^+\)。则

\[
\|\widehat T-T\|_F\le
\delta_F=\frac{e_1+\|\widehat T\|_2e_0}{s-e_0}.
\]

\end{minipage}\par

\par
\noindent\begin{minipage}{\linewidth}

\textbf{证明。} 满行秩给出

\[
\widehat T-T=(E_1-TE_0)\widehat X_0^+.
\]

\end{minipage}\par

令 \(d=\|\widehat T-T\|_F\)，则 \(d\le(e_1+\|T\|_2e_0)/s\)，同时
\(\|T\|_2\le\|\widehat T\|_2+d\)。整理得到
\(d(s-e_0)\le e_1+\|\widehat T\|_2e_0\)。证毕。无需假设
\(\widehat X_1-\widehat T\widehat X_0=0\)。

\par
\noindent\begin{minipage}{\linewidth}

对预定标记 \(W\)，有

\[
|\langle\widehat T-T,W\rangle_F|\le\delta_F\|W\|_F.
\]

\end{minipage}\par

因此 \(\omega\) 的误差半径为
\(\delta_F/\|K\|_F\)，八循环归一化标记半径为 \(\delta_F/\sqrt8\)，T20 的
\(\psi_h\) 半径为
\(\delta_F\sqrt{2/(k-h)}\)。仅当合法误差区间排除竞争符号才声明方向；秩不足不是零方向。

该结论以真实扰动满足界为条件；只有输入误差界具有相应联合覆盖时才能继承统计解释。截断伪逆、参考漂移、任务转移不确定性和系统性模型失配不会自动包含在
\(e_i\) 内。浮点容差或条件数保护也不等于经过向外舍入验证的机器证明。

\Needspace{7\baselineskip}
\subsection{A9.
连续族分离、联合误差与拒判}\label{a9.-ux8fdeux7eedux65cfux5206ux79bbux8054ux5408ux8befux5deeux4e0eux62d2ux5224}

\Needspace{6\baselineskip}
\subsubsection{命题
T25：全局最坏情况距离与数据局部排除}\label{ux547dux9898-t25ux5168ux5c40ux6700ux574fux60c5ux51b5ux8dddux79bbux4e0eux6570ux636eux5c40ux90e8ux6392ux9664}

\par
\noindent\begin{minipage}{\linewidth}

在预定观测、尺度和参数范围下，令 \(\mathcal M_F,\mathcal M_B\)
为两类非空紧致预测集合，并定义

\[
d_{FB}=\inf_{f\in\mathcal M_F,\;b\in\mathcal M_B}\|f-b\|.
\]

\end{minipage}\par

\par
\noindent\begin{minipage}{\linewidth}

若两类允许的扰动是完整闭范数球，半径分别为
\(\varepsilon_F,\varepsilon_B\)，则两误差邻域不相交当且仅当

\[
d_{FB}>\varepsilon_F+\varepsilon_B.
\]

\end{minipage}\par

\textbf{证明。}
若同一点同时位于两误差邻域，三角不等式给出一对预测的距离不超过半径之和。反之，紧致性保证最小距离由
\(f,b\) 取得；若该距离不超过半径之和，在线段 \(fb\)
上可选一点同时距两端不超过各自半径，故邻域相交。证毕。非紧集下，下确界未取得时的等号需要另外处理。

这个数是设计的最坏情况承诺。若有参数可趋向零折扣、零方向或零幅度，\(d_{FB}\)
可能为零，但某一真实观测邻域仍可能完全排除一个竞争族。局部推断应检查同一个有效不确定性区域
\(\mathcal C\) 与各族的交，而不要求两个整个族处处正分离。若参考 \(B\)
有误差，将 \(\widehat B\) 直接代入条件距离也不会自动产生覆盖保证。

\textbf{算法相关的充分条件。} 若真实点到错误族距离为
\(g\)，实际观测误差为 \(e\)，最近整个族的正确判决由 \(g>2\|e\|\)
保证。若正确类只使用训练后冻结的预测，且其距真值为 \(h\)，则充分条件变为
\(g>2\|e\|+h\)。

证明是：错误族到观测至少 \(g-\|e\|\)，正确真点到观测最多
\(\|e\|\)，冻结正确点则最多
\(\|e\|+h\)。所以不能用整族距离保证为一个有迁移误差的冻结点判决背书。以上均是确定性几何，不是额外创造的抽样概率。

\Needspace{6\baselineskip}
\subsubsection{定理
T26：候选均值残差反演与单一属性声明的错误控制}\label{ux5b9aux7406-t26ux5019ux9009ux5747ux503cux6b8bux5deeux53cdux6f14ux4e0eux5355ux4e00ux5c5eux6027ux58f0ux660eux7684ux9519ux8befux63a7ux5236}

\par
\noindent\begin{minipage}{\linewidth}

给定独立训练资料后，固定状态标签、任务块、通道处理、模板、ROI、时间窗和完整候选空间。设正式观测由
\(n\) 个完整配对向量

\[
Z_i=(B_i^\top,Y_i^\top)^\top
\]

\end{minipage}\par

\par
\noindent\begin{minipage}{\linewidth}

构成，并在条件训练分布下独立同分布高斯。对于任意预定候选
\(\theta\)，定义 \(q\) 维残差

\[
E_i(\theta)=Y_i-T_\theta B_i.
\]

\end{minipage}\par

\par
\noindent\begin{minipage}{\linewidth}

设真实 \(\theta_0\) 属于候选空间，满足
\(\mathbb E Y_i=T_{\theta_0}\mathbb E B_i\)，其残差协方差正定，且
\(n>q\)。令 \(S_{E,\theta}\) 为无偏样本协方差，

\[
R(\theta)=\bar E_\theta^\top S_{E,\theta}^{-1}\bar E_\theta,
\qquad
K_q(\alpha)=\frac{q(n-1)}{n(n-q)}F_{q,n-q}(1-\alpha).
\]

\end{minipage}\par

\par
\noindent\begin{minipage}{\linewidth}

将全部连续候选反演为

\[
\widehat\Theta_\alpha
=\{\theta:R(\theta)\le K_q(\alpha)\}.
\]

\end{minipage}\par

\par
\noindent\begin{minipage}{\linewidth}

则真实参数的保留概率为
\(1-\alpha\)。若只在该集合非空且所有候选的预定属性 \(\ell(\theta)\)
均相同时声明该属性，则

\[
\Pr_{\theta_0}\{\text{作出声明且属性错误}\mid D_{\rm train}\}
\le\alpha.
\]

\end{minipage}\par

\par
\noindent\begin{minipage}{\linewidth}

\textbf{证明。} 真实残差独立同分布高斯，均值为零。经典 Hotelling
均值分布给出

\[
\frac{n(n-q)}{q(n-1)}R(\theta_0)\sim F_{q,n-q}.
\]

\end{minipage}\par

故
\(\Pr(\theta_0\in\widehat\Theta_\alpha\mid D_{\rm train})=1-\alpha\)。若真实参数保留，所有相容候选共有的属性必与真实属性相同。因此错误声明事件包含于真实参数被排除事件，其概率不超过
\(\alpha\)。证毕。

这一统计分布是经典结果；本文使用它把参考和目标的共同误差接到 SR
连续族与机制标签上。对 FF、BB、FB、BF
应以带标签的不交并参数空间反演；不同标签即使映为同一
\(T\)，也必须全部保留。证明不要求按候选族数或无穷参数点数再做
Bonferroni，因为控制的是一个固定真参数被误排除的事件。事后选择窗口、ROI、模板、状态重排或参考来源属于另一类选择，仍须独立冻结或适当调整。

该概率不是
\(\Pr(\text{错误}\mid\text{已经声明})\)，不是候选的后验概率，也不是同时覆盖全部观测等价参数的保证。它以真实机制和共同参考关系属于候选空间为条件；模型遗漏或系统漂移不能由名义
\(p\) 值修复。

\par
\noindent\begin{minipage}{\linewidth}

\textbf{参考—目标共同协方差。} 单模板时 \(B_i,Y_i\in\mathbb R^2\)。若

\[
S_Z=\begin{pmatrix}S_B&S_{BY}\\S_{YB}&S_Y\end{pmatrix},
\]

\end{minipage}\par

\par
\noindent\begin{minipage}{\linewidth}

则

\[
\bar E_T=\bar Y-T\bar B,\qquad
V_T=S_Y+TS_BT^\top-TS_{BY}-S_{YB}T^\top.
\]

\end{minipage}\par

这正是残差的样本协方差；均值残差协方差另除以
\(n\)。配对内相关完全允许，配对之间相关不由该定理覆盖。

\textbf{多个模板。} 对 \(r\)
个固定二维模板，将所有参考、目标保留在同一个实际重复单位中。共同算子残差通常维数为
\(q=2r\)，使用 \(I_r\otimes T_\theta\) 和完整交叉协方差，阈值为
\(K_{2r}\)。模板不是额外独立被试。已知确定冗余可预先消去；不能任意截断评价协方差再沿用未经推导的普通
\(F\) 阈值。

\Needspace{6\baselineskip}
\subsubsection{命题
T27：二维残差与四维联合均值椭球不能混同}\label{ux547dux9898-t27ux4e8cux7ef4ux6b8bux5deeux4e0eux56dbux7ef4ux8054ux5408ux5747ux503cux692dux7403ux4e0dux80fdux6df7ux540c}

\par
\noindent\begin{minipage}{\linewidth}

单模板且 \(S_Z\succ0\) 时，令 \(L_T=(-T,I)\)。则

\[
\inf_{\mu:\,\mu_Y=T\mu_B}
(\bar Z-\mu)^\top S_Z^{-1}(\bar Z-\mu)
=
(\bar Y-T\bar B)^\top
(L_TS_ZL_T^\top)^{-1}
(\bar Y-T\bar B)
=R(T).
\]

\end{minipage}\par

\par
\noindent\begin{minipage}{\linewidth}

\textbf{证明。} 令 \(h=\bar Z-\mu\)，约束为
\(L_Th=L_T\bar Z\)。对严格凸二次型作拉格朗日优化，最小化向量为

\[
h_*=S_ZL_T^\top(L_TS_ZL_T^\top)^{-1}L_T\bar Z.
\]

\end{minipage}\par

代回得到等式。证毕。

\par
\noindent\begin{minipage}{\linewidth}

因此二维残差阈值

\[
K_2(\alpha)=\frac{2(n-1)}{n(n-2)}F_{2,n-2}(1-\alpha)
\]

\end{minipage}\par

\par
\noindent\begin{minipage}{\linewidth}

可以通过一个四维\textbf{辅助}椭球的相交计算实现，但该辅助椭球不是四维均值的
\(1-\alpha\) 置信域。真正的联合均值置信域是

\[
\mathcal C_{1-\alpha}
=\{\mu:(\bar Z-\mu)^\top S_Z^{-1}(\bar Z-\mu)\le K_4(\alpha)\},
\]

\end{minipage}\par

\[
K_4(\alpha)=\frac{4(n-1)}{n(n-4)}F_{4,n-4}(1-\alpha),
\qquad n>4.
\]

其覆盖率由四维 Hotelling 分布保证。将同一个 \(\mathcal C\)
与所有候选均值关系集合相交，可同时保护由真实均值关系推出的性质；任意额外均值性质或参考漂移范围若要沿用这条区域路线，必须使用
\(K_4\) 或另一个真实有效的联合区域。

二维候选残差理论只需 \(n>2\)
及残差协方差可逆；采用四维正定辅助几何的具体实现需
\(n>4\)，不可混写。作为纯阈值算例，\(n=17,\alpha=0.05\) 时
\(K_2\approx0.4621\)、\(K_4\approx0.9206\)；该算例不认证任何实际 17
个配对的独立性。

\Needspace{6\baselineskip}
\subsubsection{命题
T28：有向面积偏差与有界参考漂移敏感性}\label{ux547dux9898-t28ux6709ux5411ux9762ux79efux504fux5deeux4e0eux6709ux754cux53c2ux8003ux6f02ux79fbux654fux611fux6027}

\par
\noindent\begin{minipage}{\linewidth}

设单次配对
\(B_i=b+\varepsilon_{B,i}\)、\(Y_i=y+\varepsilon_{Y,i}\)，噪声均值为零，交叉协方差为
\(\Sigma_{BY}\)。则

\[
\mathbb E\det(B_i,Y_i)
=\det(b,y)+(\Sigma_{BY})_{12}-(\Sigma_{BY})_{21}.
\]

\end{minipage}\par

\par
\noindent\begin{minipage}{\linewidth}

若 \(n\) 次配对 iid，

\[
\mathbb E\det(\bar B,\bar Y)
=\det(b,y)+\frac{(\Sigma_{BY})_{12}-(\Sigma_{BY})_{21}}n.
\]

\end{minipage}\par

\textbf{证明。}
展开行列式，线性噪声项期望为零，交叉二次项分别为两个交叉协方差元素；样本均值的交叉协方差再除以
\(n\)。证毕。因此逐次面积平均或其 plug-in 均值不自动是无偏机制指标；T26
直接检验总体均值残差，避免把面积乘积当作无偏的简单单变量量。

\par
\noindent\begin{minipage}{\linewidth}

若实际目标参考为 \(b+\delta\)，外部依据保证 \(\|\delta\|\le\rho\)，则

\[
\inf_{\|\delta\|\le\rho}\det[b+\delta,y]
=\det[b,y]-\rho\|y\|.
\]

\end{minipage}\par

\par
\noindent\begin{minipage}{\linewidth}

对 \(y\ne0\)，将方向推至零边界的最小参考变化为

\[
\rho_*=\frac{|\det[b,y]|}{\|y\|}.
\]

\end{minipage}\par

\textbf{证明。}
\(\det[\delta,y]=\delta^\top(y_2,-y_1)^\top\)，其系数范数为
\(\|y\|\)。Cauchy–Schwarz 给极值，取 \(\delta\)
与系数向量反向即可达到；解零行列式的最小范数线性约束得到第二式。证毕。

\par
\noindent\begin{minipage}{\linewidth}

若采样误差也存在，可在真正联合均值置信域 \(\mathcal C\) 上检验

\[
\inf_{(b,y)\in\mathcal C}
\{\det[b,y]-\rho\|y\|\}>0.
\]

\end{minipage}\par

当区域覆盖和外部漂移界同时有效时，该条件保证允许漂移下的正方向。若
\(\rho\)
只是人为扫描的一组数值，输出只是敏感性分析；它不能证明实际漂移已被涵盖，也不能用
T27 的 \(K_2\) 辅助椭球冒充联合域。

\Needspace{7\baselineskip}
\subsection{A10.
论文中每种结论的正确用途}\label{a10.-ux8bbaux6587ux4e2dux6bcfux79cdux7ed3ux8bbaux7684ux6b63ux786eux7528ux9014}

\begingroup
\small

\begin{longtable}[]{@{}
  >{\raggedright\arraybackslash}p{(\columnwidth - 4\tabcolsep) * \real{0.3333}}
  >{\raggedright\arraybackslash}p{(\columnwidth - 4\tabcolsep) * \real{0.3333}}
  >{\raggedright\arraybackslash}p{(\columnwidth - 4\tabcolsep) * \real{0.3333}}@{}}
\toprule\noalign{}
\begin{minipage}[b]{\linewidth}\raggedright
要回答的问题
\end{minipage} & \begin{minipage}[b]{\linewidth}\raggedright
核心结果
\end{minipage} & \begin{minipage}[b]{\linewidth}\raggedright
实施中不可省略的条件
\end{minipage} \\
\midrule\noalign{}
\endhead
\bottomrule\noalign{}
\endlastfoot
规范距离能否决定前后向？ & T02–T03 &
明确距离、度量、状态集合与读出；同折扣比较 \\
完整相对算子拟合能否自动解除参考歧义？ & T04、T07 &
已学习参考和自由当前项必须进入竞争族 \\
一个方向属性是否由数据确定？ & T05–T06、T13–T14 &
所有合法相容解释共享该属性；保留切换和混合标签 \\
什么互反实验支持方向判断？ & T08–T09、T11 &
共同映射、已知任务结构及相应模型范围 \\
有多少观测秩才足够？ & T10、T12 &
区分完整算子、指定功能量与受限性质；检查真实投影可见 \\
定性方向后还可量化什么？ & T15、T17–T21 &
尖锐参数区间，或三个传播形状、双平面、有限核约束 \\
单平面能否证明单一 SR 核？ & T06、T16 &
不能；额外模板只检查共同作用，核识别需额外任务信息 \\
时间持续性会不会制造后瞻模式？ & T23 &
记录历史替代机制与实际探测依赖需独立检查 \\
有误差时何时可以声明？ & T24–T28 &
有效误差界或合法抽样覆盖；连续族完整排除及数值可靠性 \\
\end{longtable}

\endgroup

最终推断区分四种输出：非空相容集合具有唯一预定属性；多个属性仍相容；所有规定模型均不相容；数据或推断条件不足。最后一类包括重复单位不足、协方差不可用、相关信号不可见或覆盖假设无依据。在
T06
的单平面自由均值并集中，均值饱和性通常使“所有模型均不相容”无法出现；额外任务块和共享参数限制才提供真正可失败的核适配检查。

实际独立重复应按完整配对图谱的真实抽样单位定义，而不是按重组 occurrence
数、读出模板数或施加的计算变换数定义。不重复事件编号不等于没有共享时间支持；共同预处理、慢漂移和长时程相关均需审计。跨参与者也不能在个体参考与机制不同的情况下直接当作同一机制的
iid 重复。

\Needspace{7\baselineskip}
\subsection{A11.
代数复现与结果索引}\label{a11.-ux4ee3ux6570ux590dux73b0ux4e0eux7ed3ux679cux7d22ux5f15}

\path{code/math_checks.py} 与 \path{code/math_additions.py}
复现精确构造、补偿公式、逆解和扰动实例。计算输出与部分折扣界位于
\path{source_data/}；登记表 \path{math/theorem_registry.json} 索引
T01–T28。一般结论由以上证明给出，真实测量方法及来源见 B 部分。

\clearpage

\section{B.
真实重复测量与校准}\label{b.-ux771fux5b9eux91cdux590dux6d4bux91cfux4e0eux6821ux51c6}

\Needspace{7\baselineskip}
\subsection{E1.
目的与证据层级}\label{e1.-ux76eeux7684ux4e0eux8bc1ux636eux5c42ux7ea7}

经验分析把主文的测量方法接到实际神经记录上：参考模式能否预测较晚的测量，任务相容读出能否在已知变换下保留方向信息，采集结构与参考迁移又怎样影响判决。已知变换校准把指定算子作用于实测反应，使变换本身具有已知真值，同时保留测量中的实际波动。

三个层级分别报告：①实际较早/较晚测量上的参考预测；②给定配对误差假设的已知变换校准；③原始采集层级和跨记录迁移的限制。预测成功、条件性方向排除及独有
SR 生物机制证明并不是同一终点。

\Needspace{7\baselineskip}
\subsection{E2.
数据、状态和采集单位}\label{e2.-ux6570ux636eux72b6ux6001ux548cux91c7ux96c6ux5355ux4f4d}

主图采用 Turner et al.~(2025) 的圆周位置 EEG 定位数据子集
{[}E1{]}。原论文的参与者数为
18，本分析仅用先前下载的一位参与者的两次定位记录；旧包将其标记为
ID13。每次记录的供应数组含 1980 个合格事件、40
个位置标签和两个实际块。两个块来自定位过程中的休息，不是两名参与者。选用与原论文相同的
17
个枕顶电极：P7、P5、P3、P1、Pz、P2、P4、P6、P8、PO7、PO3、POz、PO4、PO8、O1、Oz、O2。

早窗为刺激后 100–150 ms，晚窗为 200–250 ms；两者减去 \(-200\)–0 ms
的基线。分析输入为已归档校准包中的事件级派生测量。NPZ 以
\texttt{early}、\texttt{late} 标记时窗；精确时窗范围和 ID13
的文件对应关系记录于来源统计附录，而没有作为提取索引或参与者字段嵌入
NPZ。因此，本校准的可复现起点是所附事件级测量数组。

跨记录压力分析另使用 Kern MEG 的已有定位重复，公开档案 DOI 为
10.5281/zenodo.8001755 {[}E2{]}。Ekman et al.~定位器的两个参与者、V1
和海马四个情境仅作补充参考估计例
{[}E3{]}。本文不把这些所选子集当成相应原研究的完整参与者样本。

\Needspace{7\baselineskip}
\subsection{E3.
参考估计：训练与评价分开，所有合格控制事件可用}\label{e3.-ux53c2ux8003ux4f30ux8ba1ux8badux7ec3ux4e0eux8bc4ux4ef7ux5206ux5f00ux6240ux6709ux5408ux683cux63a7ux5236ux4e8bux4ef6ux53efux7528}

\Needspace{6\baselineskip}
\subsubsection{E3.1
时间划分与估计候选}\label{e3.1-ux65f6ux95f4ux5212ux5206ux4e0eux4f30ux8ba1ux5019ux9009}

在每个 Turner
实际块内，按事件起始时间的中点划分较早控制段与较晚评价段，并在中点两侧各排除
0.5 s。8 个“记录\(\times\)块\(\times\)时窗”情境的控制事件数为 491 或
492；评价事件数为 492 或 495。不同状态的控制重复数允许不等，范围为每状态
5–22
个；不为凑平衡而丢弃其他可用控制事件。训练和评价事件编号不交叠，边界最近事件起始间隔约
1.20 s。该间隔避免直接边界复用，不保证整段 EEG 误差相互独立。

控制内交叉验证为每个状态留出一个事件；每一轮只用其他控制事件的状态均值拟合。候选包括直接均值、零模式、标量收缩、截断奇异值估计，以及保持正弦/余弦成对的
Fourier 状态平滑。收缩系数在 0–1、步长 0.1 的集合中选择。Fourier
截止候选保留 1、2、4、8、19
个谐波对或完整中心化状态空间。成对保留谐波不会单独偏爱前向或后向旋转。候选的选择依据为控制内完整矩阵
MSE，不使用外层后段或方向真值标签。

主文经验图显示 \path{cv_fourier}
候选族的外层结果，每个点对应一个实际块中的较早训练和较晚评价。估计方法的开发与评价属于回顾性分析；在每次拟合内部，后段仅供评价，参数选择全部使用较早控制段。

\Needspace{6\baselineskip}
\subsubsection{E3.2
任务相关平面上的损失与零基线}\label{e3.2-ux4efbux52a1ux76f8ux5173ux5e73ux9762ux4e0aux7684ux635fux5931ux4e0eux96f6ux57faux7ebf}

设状态中心化参考估计为
\(\widehat B\in\mathbb R^{40\times17}\)，后段状态均值为
\(X\)。固定第一圆周谐波平面为

\[
C=\sqrt{\frac{2}{40}}
\begin{bmatrix}
\cos(2\pi j/40)\\
\sin(2\pi j/40)
\end{bmatrix}_{j=0}^{39}.
\]

任务平面上的两项损失为

\[
L_C(\widehat B,X)=\frac{\|C(\widehat B-X)\|_F^2}{2\times17},
\qquad L_C(0,X)=\frac{\|CX\|_F^2}{2\times17}.
\]

单块收益为
\(G_C=1-L_C(\widehat B,X)/L_C(0,X)\)。零矩阵是在已经去除所有状态共同平均活动之后的基线，表示未预测任何状态差异；它不表示脑没有活动。采用零基线还可识别一种重要情况：一个强收缩估计可以大幅击败噪声很大的直接均值，同时不保留任何可用于辨别状态的信号。

若同一后段满足 \(X=B_*+\varepsilon\) 且
\(E[\varepsilon\mid\text{训练资料}]=0\)，则两个固定训练估计的期望预测损失差等于它们对
\(B_*\) 的平方误差差，测试噪声项抵消。但一次实际后段损失不是 \(B_*\)
的确定误差，也不是误差上界。跨段漂移和条件均值变化会改变这一解释。

\Needspace{6\baselineskip}
\subsubsection{E3.3 结果}\label{e3.3-ux7ed3ux679c}

\Needspace{9\baselineskip}

\begingroup
\small

\begin{longtable}[]{@{}lrr@{}}
\toprule\noalign{}
实际记录与块 & 早窗：相对零模式收益 & 晚窗：相对零模式收益 \\
\midrule\noalign{}
\endhead
\bottomrule\noalign{}
\endlastfoot
记录 1，块 1 & \(-8.3292\)\% & 16.9238\% \\
记录 1，块 2 & 0.0000\% & 6.8747\% \\
记录 2，块 1 & 25.1302\% & 32.3970\% \\
记录 2，块 2 & 29.3631\% & 57.9619\% \\
\end{longtable}

\endgroup

按四块损失先平均再取比值，早窗的平面 MSE 从零模式的 12.094698 降至
10.313989，收益 \textbf{14.7231\%}；晚窗从 5.998847 降至 4.386882，收益
\textbf{26.8713\%}。这是两个时窗的分开评价，不合并为独立样本。

晚窗四块均改善；早窗仅第二次记录的两个块改善。因此已有证据支持这个子集中特定平面的可预测性，但不支持“同一人在各次记录均稳定”的一般断言。更进一步，晚窗完整矩阵平均
MSE
仍高于完整矩阵零基线；可预测低维平面不等于恢复全部状态算子。参考预测最好的时窗也不必是方向排除最有力的时窗，两者针对不同损失方向。

这些块间比较描述实际后段的预测损失。潜在参考精度是否稳定、各块的期望预测风险是否相等，属于另外的统计问题。

\Needspace{7\baselineskip}
\subsection{E4. 已知 SR
变换校准和任务相容读出}\label{e4.-ux5df2ux77e5-sr-ux53d8ux6362ux6821ux51c6ux548cux4efbux52a1ux76f8ux5bb9ux8bfbux51fa}

\Needspace{6\baselineskip}
\subsubsection{E4.1
重组状态图与训练限制}\label{e4.1-ux91cdux7ec4ux72b6ux6001ux56feux4e0eux8badux7ec3ux9650ux5236}

为获得重复的完整状态矩阵，按每个状态第 k 次出现拼成第 k 张状态图。Turner
第二次记录各状态的最小出现次数为 38，故形成 38 张 40(\times)17
的矩阵。前 4 张仅用于拟合参考、通道白化与读出；其余 34 张交替构成 17
对参考/目标基础测量。前者使用 160 个原始事件，后者使用 1360
个，编号不交叠。

这种拼接没有创造新的独立采集。38 张状态图中 18
张横跨两个实际块，单张的时间跨度中位数为 133.318 s。按“第 k
次出现”聚合的时间区间相互重叠，尽管不是同一事件被重复使用。

训练参考估计在标量、SVD、Fourier
三族间用控制内损失选择。选定参考后可保持原生状态次序，或只用训练参考的二维主成分，通过确定性符号固定与线性指派安排圆周状态顺序。后者使四分之一周期平面
\(k=10\)
捕获较强的训练对比。由此得到的是可供未来实验实现的任务顺序；现有参与者没有在校准步骤中重新学习该任务。

\Needspace{6\baselineskip}
\subsubsection{E4.2
读出保持算子作用}\label{e4.2-ux8bfbux51faux4fddux6301ux7b97ux5b50ux4f5cux7528}

令 \(Z=C_{\rm eff}XW\)
为固定状态平面与固定通道白化后的测量，\(R=C_{\rm eff}\widehat BW\)
为训练模板，\(J=\left[\begin{smallmatrix}0&-1\\1&0\end{smallmatrix}\right]\)。采用

\[
\mathcal L_R(Z)=\frac{1}{\|R\|_F^2}
\begin{bmatrix}\langle R,Z\rangle_F\\\langle JR,Z\rangle_F\end{bmatrix}.
\]

\textbf{实现与理论的归一化。} 理论附录 T12 中的读出使用
\(1/\|R\|_F\)，经验实现使用上式的 \(1/\|R\|_F^2\)。令
\(s=1/\|R\|_F>0\)，则实现读出等于理论读出再乘训练后冻结的共同正尺度
\(s\)；参考与目标同时缩放。因而两向量的相对角不变，有向面积乘 \(s^2\)
而符号不变，任务交换关系也不变。对同一个候选算子，残差及其样本协方差分别变为
\(r_T'=s r_T\)、\(V_T'=s^2V_T\)，所以在非退化协方差和精确算术下

\[
(r_T')^\top(V_T')^{-1}r_T'=r_T^\top V_T^{-1}r_T.
\]

因此，固定样本量和同一名义阈值下，这两个归一化约定给出同一候选残差反演。该等价仅比较\textbf{其余训练结果与候选族已冻结后的共同尺度}；重新选择白化、模板、正则化强度或其他处理可能改变读出，不能由这个恒等式声称所有正则化参数或拟合结果自动不变。实际范数或协方差退化时仍需拒判。

该读出满足 \(\mathcal L_R(JZ)=J\mathcal L_R(Z)\)，因而保留平面上
\(uI+vJ\)
的作用。若模板为零或指定平面近乎消失，流程拒判，不用其他数据临时补出有利模板。白化、状态排列、模板和读出全部在原训练部分固定，随后对参考和目标采用同一映射。

\Needspace{6\baselineskip}
\subsubsection{E4.3
已知变换到底改变了什么}\label{e4.3-ux5df2ux77e5ux53d8ux6362ux5230ux5e95ux6539ux53d8ux4e86ux4ec0ux4e48}

对二维基础测量 \(v_i\) 施加已知 \(T\)，构造 \(y_i=Tv_i\)。规范 SR 使用
\(T=(I-\gamma A)^{-1}\) 或 \((I-\gamma A^\top)^{-1}\)，其中 \(A\)
为指定平面上的循环算子；校准折扣集合为
.25、.5、.8、.95。其他已保存的计算条件还包含独立当前增益及预测增益。

这里变换的是整份真实测量，信号和噪声同时变化。没有生成虚拟参与者或附加模拟噪声，但目标也不是自然学习任务新采集的神经活动。因此它能够检验映射的代数行为和真实重复噪声下的条件性可辨别程度，不能估计固定生物噪声下的通常实验功效，更不能证明
SR 是唯一生物机制。

该实证仅校准给定参考关系下的方向判据，没有让参与者实施本文提出的互反任务，也没有验证纯完整
SR 的 FF/BB/FB/BF 排他区。已知规范后向 SR
规定了校准中施加的算子。该标签对应变换真值；实际测量终点是这一变换下的方向排除。

\Needspace{7\baselineskip}
\subsection{E5.
联合误差、连续排除与三个条件性阳性}\label{e5.-ux8054ux5408ux8befux5deeux8fdeux7eedux6392ux9664ux4e0eux4e09ux4e2aux6761ux4ef6ux6027ux9633ux6027}

用每对 \((b_i,y_i)\)
的四维联合样本估计均值和协方差。主要规则是\textbf{候选算子的二维配对残差检验反演}，其辅助椭球不是四维均值的同时置信椭球。方向判决通过连续优化边界完成，并非只检查几个折扣网格点。条件性覆盖要求给定训练资料后各配对单元独立同分布高斯、共同的未变换参考均值关系正确、真实候选的残差协方差非奇异，并且训练参数冻结且与评价保持所需独立性。仅有配对可交换性或不重复事件编号并不足够。

主文经验图报告零方向边界的最小二次统计量与临界值之比。17 对样本、名义
\(\alpha=.05/54\) 下，临界值为 1.447336。54 对应档案中的 27
个测量/任务情境乘两个模板版本；每个情境还包含多个已知变换条件，因此不是整张结果表只有
54 行。比值超过 1
表示给定该误差模型时，相关方向边界被排除；线段表示超出边界的余量，不是误差条。

\Needspace{8\baselineskip}

\begingroup
\small

\begin{longtable}[]{@{}
  >{\raggedright\arraybackslash}p{(\columnwidth - 8\tabcolsep) * \real{0.1579}}
  >{\raggedleft\arraybackslash}p{(\columnwidth - 8\tabcolsep) * \real{0.2105}}
  >{\raggedleft\arraybackslash}p{(\columnwidth - 8\tabcolsep) * \real{0.2105}}
  >{\raggedleft\arraybackslash}p{(\columnwidth - 8\tabcolsep) * \real{0.2105}}
  >{\raggedleft\arraybackslash}p{(\columnwidth - 8\tabcolsep) * \real{0.2105}}@{}}
\toprule\noalign{}
\begin{minipage}[b]{\linewidth}\raggedright
固定训练读出
\end{minipage} & \begin{minipage}[b]{\linewidth}\raggedleft
已知后向 SR 折扣
\end{minipage} & \begin{minipage}[b]{\linewidth}\raggedleft
边界统计量/临界值
\end{minipage} & \begin{minipage}[b]{\linewidth}\raggedleft
名义 54 情境校正 p 值
\end{minipage} & \begin{minipage}[b]{\linewidth}\raggedleft
相位
\end{minipage} \\
\midrule\noalign{}
\endhead
\bottomrule\noalign{}
\endlastfoot
原生次序，k=1 & .95 & 1.206961 & .020613 & \(-67.4544\)° \\
设计次序，k=10 & .80 & 1.156972 & .025294 & \(-38.6598\)° \\
设计次序，k=10 & .95 & 1.341880 & .012182 & \(-43.5312\)° \\
\end{longtable}

\endgroup

三者复用同一人的第二次记录、同一个早窗和相同 17
对测量；不存在三次独立复制。原生 k=1 的原 A 与 B-trained A
结果因共同缩放而统计等价，图中只保留一份。设计平面的两个阳性使用更新训练模板；正标量缩放本身不会提高尺度不变判决的辨别力。

上述三个边界余量对应统一放大协方差约 20.70\%、15.70\%、34.19\%
后恰触及名义阈值。它们是敏感性几何，不是已估计出来的实际 EEG
依赖修正。该门槛下的前向真值均没有方向声明，故不报告对称的双向正确率。

名义校正是在配对误差模型下对指定的 54
个情境进行控制。方法属于回顾性开发序列，这一校正不覆盖整个方法选择过程。时间依赖、共享上游处理以及参考漂移，决定该条件性抽样模型是否符合一次真实采集。因此，这三项结果按条件性校准结果报告。

\Needspace{7\baselineskip}
\subsection{E6.
实际块与跨记录压力检验}\label{e6.-ux5b9eux9645ux5757ux4e0eux8de8ux8bb0ux5f55ux538bux529bux68c0ux9a8c}

在设计 k=10、标准后向 \(\gamma=.95\)
条件下，分别限定两个实际块重新形成完整状态重复，可用配对数降为 6 与
7；相同名义联合门槛下两块均拒判。既改变了有效重复数，也改变了均值及协方差，不能把结果唯一归因于“时间相关造成假阳性”。若坚持仅以两个原始块作为独立单位，又不足以估计所用四维联合协方差。这正是后续任务应增加独立采集的设计约束。

跨记录压力分析保持来源训练得到的 \(C_{\rm eff}\)、\(W\)、\(R\) 与
\(H\)，不以目标记录重新拟合有利坐标。同记录和跨记录比较使用相同来源参考与匹配的目标样本预算。名义联合门槛下，\textbf{1200/1200}
个跨记录计算条件拒判，其中 Turner 570、Kern
630。不同折扣、变换族和读出多次重用同一记录，这些数字既不是参与者数，也不是独立试验次数。

跨记录可能改变真正参考均值；因此该压力试验不是已知零假设下的错误率估计。全部拒判不证明参考相等，也不证明脑没有
SR，只说明当前流程未在这些参考迁移条件下给出足够支持的唯一方向声明。

\Needspace{7\baselineskip}
\subsection{E7. 补充 fMRI
参考示例及其边界}\label{e7.-ux8865ux5145-fmri-ux53c2ux8003ux793aux4f8bux53caux5176ux8fb9ux754c}

Ekman 数据使用已处理的定位时序和提取出的模式矩阵。每位参与者、每个 ROI
有 8 个位置(\times)8 次事件；每状态前 4 次训练、后 4
次评价。真实事件起始时间与已知提取窗口支持该时间划分，但上游全记录处理及血流动力学相关仍然存在。

S01 V1 的控制内选择给出 \(\lambda\)=.3；完整模式外层 MSE 为
.05297993，零模式为 .06589265，直接均值为 .06411860；分别改善 19.5966\%
和约 17.37\%。S01/S02 的海马控制内估计均选零；S02 V1 的标量/SVD
也选零，Fourier 的零基线收益约 .5454\%。全部情境保留于
\path{source_data/empirical_ekman_reference_all.csv}。这个补充例子说明参考估计必须与零基线比较，不能据此诊断该研究的方向错误或建立普遍定位器污染幅度。

\Needspace{7\baselineskip}
\subsection{E8.
数值核对与可复现的分析单位}\label{e8.-ux6570ux503cux6838ux5bf9ux4e0eux53efux590dux73b0ux7684ux5206ux6790ux5355ux4f4d}

经验材料保留事件级输入、来源表格，以及把它们连接到报告结果的核对程序。便携脚本
\path{code/empirical_check_fig07_portable.py} 从所附实测数组重新拟合 40
个 Turner
参考估计，重建阳性条件使用的两个冻结读出，并核对预测收益、边界比值和跨记录结果的汇总。参考矩阵最大差为
0，损失最大差为 \(3.55\times10^{-15}\)；两个冻结读出的最大差也为
0。训练事件身份以及训练、评价事件不交叠均与归档输入一致。这些计算以派生测量为起点，原始
EEG 预处理和采集层面的独立性模型位于这个可复现单位的上游。

对于较早的连续模型族校准，\path{code/empirical_check_dfb_tables.py}
联结逐折判决表与观测类别审计，验证充分不等式，重建 E9
的计数。来源文件哈希和原始归档路径保存在
\path{empirical/dFB_provenance.json}、\path{empirical/historical_source_manifest.json}
与 \path{empirical/round1_source_manifest.json}。核对结果为
\path{empirical/fig07_portable_checks.json} 和
\path{empirical/dFB_table_checks.json}。图中数值直接读取相应 CSV。

各项经验终点按其实际单位报告：参考估计对应事件级预测；已知变换校准对应计算条件；采集和迁移检查对应实际块与记录；只有以人为抽样单位的分析才按参与者进行推断。重复一种计算条件不会增加独立采集次数。

\Needspace{7\baselineskip}
\subsection{E9. Kern 与 Turner
实际重复测量中的连续族分离}\label{e9.-kern-ux4e0e-turner-ux5b9eux9645ux91cdux590dux6d4bux91cfux4e2dux7684ux8fdeux7eedux65cfux5206ux79bb}

\Needspace{6\baselineskip}
\subsubsection{E9.1
数据结构与候选模型族}\label{e9.1-ux6570ux636eux7ed3ux6784ux4e0eux5019ux9009ux6a21ux578bux65cf}

本校准采用 Kern MEG 档案中的一位参与者 DSMR101 {[}E2, E4{]}，以及 Turner
EEG 档案中的一位参与者 ID13
{[}E1{]}。每个数据源提供两个反应时窗、三个状态空间和四个复用参考划分。三个状态空间分别是单
Fourier 平面、双 Fourier 平面和完整中心化状态空间；Kern 的维数为
2、4、11，Turner 为 2、4、39。因此，共有 \(2\times2\times3\times4=48\)
个参考折设置。

每个设置包含 18 个已知方向变换：两个方向、三个折扣
\(\gamma_0\in\{.2,.5,.8\}\)、三个预测强度
\(\rho\in\{.1,.3,1\}\)，另有一个仅当前对照。逐折共 912
条条件，其中方向条件为 864 条。每四折合并后为 228 条条件，其中 216
条为方向条件。这些计数描述对两位参与者记录施加的变换与划分。

\par
\noindent\begin{minipage}{\linewidth}

拟合候选为

\[
T_j(c,\alpha,\gamma)S=cS+\alpha N_\gamma(A_j)S,
\qquad N_\gamma(A)=A(I-\gamma A)^{-1},
\]

\end{minipage}\par

其中当前系数 \(c\in\mathbb R\) 自由，预测增益 \(\alpha\geq0\)，折扣
\(\gamma\in[0,.95]\)。前向、后向候选分别使用 \(A\) 和
\(A^\top\)；对称转移候选使用
\((A+A^\top)/2\)，另设仅当前模型。循环任务上产生相同预测的两种对称公式合并为一个观测类别。若
\(\alpha\leq10^{-12}\)，或者拟合的预测在观测上等价于仅当前，则不论候选的名义方向标签为何，一律归入当前类。

\par
\noindent\begin{minipage}{\linewidth}

令 \(R,V\) 为训练参考和训练目标基础测量，\(S,Z\)
为相应测试测量。已知算子 \(T_0\) 产生目标 \(T_0V\) 和
\(T_0Z\)，信号与测量波动一起接受变换。施加的预测强度按训练目标基础输入定义：

\[
\rho=\frac{\alpha_0\|Q_VN_{\gamma_0}(A_j)V\|_F}{\|V\|_F},
\]

\end{minipage}\par

其中 \(Q_V\) 去除向量化的当前模式方向。在另一个参考 \(S\)
上，仍使用同一冻结的 \(\alpha_0\) 计算预测残差，其相对大小不必保持为
\(\rho\)。通道白化由训练参考估计。Kern
的上游信号空间投影由完整第一定位记录估计，因此后续四个划分共享这一预处理。

\Needspace{6\baselineskip}
\subsubsection{E9.2
固定真实点与连续竞争族}\label{e9.2-ux56faux5b9aux771fux5b9eux70b9ux4e0eux8fdeux7eedux7adeux4e89ux65cf}

允许
\(\alpha=0\)，意味着每个方向族都与当前族相交。因此，不加限制的前向、后向全族最小距离为零。这里有用的量是一个具有非零预测成分的指定真实预测，到整个竞争模型族的距离。

\par
\noindent\begin{minipage}{\linewidth}

在固定参考 \(S\) 上，记

\[
r_S=\alpha_0\|Q_SN_{\gamma_0}(A_j)S\|_F,
\qquad
u_0=\frac{Q_SN_{\gamma_0}(A_j)S}{\|Q_SN_{\gamma_0}(A_j)S\|_F}.
\]

\end{minipage}\par

\par
\noindent\begin{minipage}{\linewidth}

竞争族 \(k\) 的预测方向在同样消去当前项并归一化后记为
\(\nu_k(\gamma)\)。在实数域优化当前系数，并在 \([0,\infty)\)
上优化预测增益，可得

\[
m_{j\rightarrow k,S}
=r_S\sqrt{1-\left[\max\left\{0,
\sup_{\gamma\in[0,.95]}\langle\nu_0,\nu_k(\gamma)\rangle_F
\right\}\right]^2}.
\]

\end{minipage}\par

若竞争预测残差为零，该候选属于当前观测类，距离为
\(r_S\)；只对非零残差进行归一化。二元距离只比较前向和后向；多元距离还包括对称和仅当前解释，并取其中的最小值。连续优化给出点到竞争族分离距离的下界
\(\underline m_S\)。计算采用数值解析区间界，并有独立 witness
核对；这些是浮点计算界，没有使用向外舍入的区间算术。拟合采用的有限折扣网格与排除采用的连续参数域分别记录。

\Needspace{6\baselineskip}
\subsubsection{E9.3
哪一部分误差会改变模型比较}\label{e9.3-ux54eaux4e00ux90e8ux5206ux8befux5deeux4f1aux6539ux53d8ux6a21ux578bux6bd4ux8f83}

\par
\noindent\begin{minipage}{\linewidth}

一个设置中的全部候选预测均落在共同空间

\[
U_S=\operatorname{span}\{\operatorname{vec}(S),
\operatorname{vec}(AS),\ldots,\operatorname{vec}(A^{q-1}S)\}
\]

\end{minipage}\par

\par
\noindent\begin{minipage}{\linewidth}

内，其中 \(q\)
是状态对比数。对于这些循环算子，多项式约化也覆盖后向和对称候选。该空间维数至多为
\(q\)，完整测量却有 \(q\) 乘通道数个坐标。对任意候选 \(f\in U_S\)，

\[
\|Y-f\|^2=\|\Pi_{U_S}Y-f\|^2+
\|(I-\Pi_{U_S})Y\|^2.
\]

\end{minipage}\par

最后一项对所有候选相同，删除它精确保留候选间的平方误差差值以及最近族判决。

\par
\noindent\begin{minipage}{\linewidth}

校准中心为 \(\mu_S=T_0S\)，故变换后测试测量的实际偏差为

\[
E=T_0(Z-S),\qquad T_0Z=\mu_S+E.
\]

\end{minipage}\par

\par
\noindent\begin{minipage}{\linewidth}

864 个方向条件中的比值中位数为

\[
\frac{\|\Pi_{U_S}E\|}{\|E\|}=.596376,
\qquad
\frac{\|Q_S\Pi_{U_S}E\|}{\|E\|}=.191525.
\]

\end{minipage}\par

这是误差范数之比，不是解释方差百分比，也不是预测 MSE
的下降比例。第二个投影适用于所有候选都能在当前待分类目标上、于实数域自由重新拟合当前系数的规则。若当前系数在训练后冻结，判决规则便不同，须计入它自己的迁移误差。

在 48 个参考设置中，独立 SVD 与谱方法所得 \(\Pi_{U_S}\) 最大差为
\(6.01\times10^{-13}\)；4,560 个归档候选评分中的分数差恒等式最大差为
\(1.25\times10^{-12}\)。归档输入复原对训练、测试 SSE 的最大差分别为
\(3.98\times10^{-13}\) 与 \(2.28\times10^{-13}\)。

\Needspace{6\baselineskip}
\subsubsection{E9.4
与最近模型族规则相配的充分条件}\label{e9.4-ux4e0eux6700ux8fd1ux6a21ux578bux65cfux89c4ux5219ux76f8ux914dux7684ux5145ux5206ux6761ux4ef6}

在目标内最近族分类中，每个候选针对同一参考 \(S\) 和目标 \(T_0Z\)
重新拟合 \(c\)、\(\alpha\) 与 \(\gamma\)。拟合网格为 0–.95、步长 .01 的
96
个折扣，包含每一个施加的真实折扣，因此正确拟合族包含真实预测。竞争解释的排除仍采用更大的连续族。

若 \(\varepsilon\) 控制会影响分数的那部分误差，正确预测离目标至多为
\(\varepsilon\)，任一竞争预测离目标至少为
\(\underline m_S-\varepsilon\)。所以 \(\underline m_S>2\varepsilon\)
足以保证最近族判决正确。代入实际校准误差，得到下列确定性核对：

\Needspace{8\baselineskip}

\begingroup
\small

\begin{longtable}[]{@{}
  >{\raggedright\arraybackslash}p{(\columnwidth - 6\tabcolsep) * \real{0.2000}}
  >{\raggedleft\arraybackslash}p{(\columnwidth - 6\tabcolsep) * \real{0.2667}}
  >{\raggedleft\arraybackslash}p{(\columnwidth - 6\tabcolsep) * \real{0.2667}}
  >{\raggedleft\arraybackslash}p{(\columnwidth - 6\tabcolsep) * \real{0.2667}}@{}}
\toprule\noalign{}
\begin{minipage}[b]{\linewidth}\raggedright
充分条件采用的误差范数
\end{minipage} & \begin{minipage}[b]{\linewidth}\raggedleft
二元通过 / 864
\end{minipage} & \begin{minipage}[b]{\linewidth}\raggedleft
多元通过 / 864
\end{minipage} & \begin{minipage}[b]{\linewidth}\raggedleft
通过条件中相配规则的错误数
\end{minipage} \\
\midrule\noalign{}
\endhead
\bottomrule\noalign{}
\endlastfoot
完整测量，\(\|E\|\) & 0 & 0 & 0 \\
共同预测空间，\(\|\Pi_{U_S}E\|\) & 6 & 6 & 0 \\
共同空间并消去当前项，\(\|Q_S\Pi_{U_S}E\|\) & 80 & 77 & 0 \\
\end{longtable}

\endgroup

通过条件中的最小正余量约为
\(6.74\times10^{-4}\)，远高于观察到的投影数值误差。这里检验的是在已知变换、误差可计算时，一个确定性充分条件是否兑现；它是事后核对，没有名义置信水平，与
E5 的 17 对残差推断属于不同终点。

同一批条件还说明，保证必须指明采用什么判决规则。较早的规则在 \(R,T_0V\)
上估计参数，随后把参数冻结，用于 \(S\)
上的预测。两种规则统一采用同一观测类别后，结果为：

\Needspace{7\baselineskip}

\begingroup
\small

\begin{longtable}[]{@{}
  >{\raggedright\arraybackslash}p{(\columnwidth - 8\tabcolsep) * \real{0.1579}}
  >{\raggedleft\arraybackslash}p{(\columnwidth - 8\tabcolsep) * \real{0.2105}}
  >{\raggedleft\arraybackslash}p{(\columnwidth - 8\tabcolsep) * \real{0.2105}}
  >{\raggedleft\arraybackslash}p{(\columnwidth - 8\tabcolsep) * \real{0.2105}}
  >{\raggedleft\arraybackslash}p{(\columnwidth - 8\tabcolsep) * \real{0.2105}}@{}}
\toprule\noalign{}
\begin{minipage}[b]{\linewidth}\raggedright
满足目标内充分条件的条件集
\end{minipage} & \begin{minipage}[b]{\linewidth}\raggedleft
冻结规则：可见方向正确
\end{minipage} & \begin{minipage}[b]{\linewidth}\raggedleft
冻结规则：仅当前
\end{minipage} & \begin{minipage}[b]{\linewidth}\raggedleft
冻结规则：对称
\end{minipage} & \begin{minipage}[b]{\linewidth}\raggedleft
目标内规则：可见方向正确
\end{minipage} \\
\midrule\noalign{}
\endhead
\bottomrule\noalign{}
\endlastfoot
二元，80 条 & 42 & 38 & 0 & 80 \\
多元，77 条 & 39 & 37 & 1 & 77 \\
\end{longtable}

\endgroup

第二行的 37 个仅当前与 1 个对称结果，都没有恢复唯一方向，不能写成 38
次前后方向反转。两种方法回答不同问题：一种检验训练参数的迁移预测，另一种在拟合当前目标后归类。因此，这项对照展示的是保证式中需要补入的误差项，而不是独立预测准确率的提高。

\par
\noindent\begin{minipage}{\linewidth}

具体地，令 \(\widehat T_{\rm true}\)
表示正确标签模型族在训练中得到的估计，则

\[
q_S=\|\widehat T_{\rm true}S-T_0S\|
\]

\end{minipage}\par

\par
\noindent\begin{minipage}{\linewidth}

测量参数迁移误差。冻结规则的一项充分条件为

\[
\underline m_S>2\|\Pi_{U_S}E\|+q_S.
\]

\end{minipage}\par

\par
\noindent\begin{minipage}{\linewidth}

没有方向条件通过这一保守界。较紧的实际残差界

\[
\underline m_S>\|\Pi_{U_S}E\|
+\|\Pi_{U_S}(T_0Z-\widehat T_{\rm true}S)\|
\]

\end{minipage}\par

分别通过 10 个二元和 10
个多元逐折条件，全部判别正确。四折合并后，没有条件通过上述任一冻结规则界。一次判决可以碰巧正确而没有通过充分条件；充分条件的作用，是圈出能由指定误差比较保证正确的区域。

\Needspace{6\baselineskip}
\subsubsection{E9.5
预测未知误差与分辨折扣}\label{e9.5-ux9884ux6d4bux672aux77e5ux8befux5deeux4e0eux5206ux8fa8ux6298ux6263}

未来目标的 \(T_0(Z-S)\)
尚不可见。因此，归档分析还用训练配对差与参考重复差构造经验误差尺度。这些量描述已经观察到的波动，不是下一份目标误差的置信上界。它们对冻结规则成功条件的排序能力因数据源而异：多元比较中，配对误差加迁移误差的指标在
Kern 的 AUROC 为 .537，在 Turner 为
.901。这里排序的是反复使用同一记录构造的计算条件；每个数据源只有一人，参与者与数据源完全混杂。

跨数据源阈值探索先在来源数据中寻找“至少 10 条条件通过、观察正确率至少
90\%”的阈值，再原样应用于另一个数据源。20
个“方向比较\(\times\)指标\(\times\)迁移方向”组合中只有两个能给出这样的来源阈值；其正确条件数在迁移后分别从
12/12 变成 9/14、从 9/10 变成
27/36，均未保持来源标准。表中保留算术并列以及排序的数值界；这些界反映优化精度，不是抽样置信区间。另一项
THINGS-fMRI
校准使用互不重叠的两组六次记录，分别用于构建和评价阈值，同样没有得到指定的可迁移低错误阈值。该项研究列于
E11 的 D05。

折扣分辨还需要回答另一个问题：真实点能否与同方向中所有
\(|\gamma-\gamma_0|\geq.2\)
的候选分开？在单个旋转平面上，删除自由当前项后只剩一个预测方向；自由预测增益可以吸收
\(\gamma\) 的影响，因此该距离精确为零。双平面则可以产生非零的形状差异：

\Needspace{11\baselineskip}

\begingroup
\small

\begin{longtable}[]{@{}
  >{\raggedright\arraybackslash}p{(\columnwidth - 6\tabcolsep) * \real{0.2000}}
  >{\raggedleft\arraybackslash}p{(\columnwidth - 6\tabcolsep) * \real{0.2667}}
  >{\raggedleft\arraybackslash}p{(\columnwidth - 6\tabcolsep) * \real{0.2667}}
  >{\raggedleft\arraybackslash}p{(\columnwidth - 6\tabcolsep) * \real{0.2667}}@{}}
\toprule\noalign{}
\begin{minipage}[b]{\linewidth}\raggedright
数据源与状态空间
\end{minipage} & \begin{minipage}[b]{\linewidth}\raggedleft
单位折扣距离下界中位数
\end{minipage} & \begin{minipage}[b]{\linewidth}\raggedleft
冻结拟合中可报告 \(\gamma\) / 144
\end{minipage} & \begin{minipage}[b]{\linewidth}\raggedleft
可报告且 \(\lvert\widehat\gamma-\gamma_0\rvert<.2\) / 144
\end{minipage} \\
\midrule\noalign{}
\endhead
\bottomrule\noalign{}
\endlastfoot
Kern，单平面 & 0 & 0 & 0 \\
Kern，双平面 & .234863 & 124 & 52 \\
Kern，完整 11 维空间 & .228635 & 120 & 43 \\
Turner，单平面 & 0 & 0 & 0 \\
Turner，双平面 & .081264 & 112 & 27 \\
Turner，完整 39 维空间 & .234912 & 144 & 71 \\
\end{longtable}

\endgroup

单平面约 \(2.11\times10^{-8}\)
的数值上界反映精确零附近的舍入。多平面的正距离建立了几何上的分辨可能性，但在这些测量中，没有折扣条件通过实际误差充分界。因此，可报告一个拟合值与保证一定的折扣分辨率是不同终点。E9
来源文件统一使用 \path{source_data/dFB_} 前缀。

\Needspace{7\baselineskip}
\subsection{E10.
参考内容、物理历史与时间权重的解释}\label{e10.-ux53c2ux8003ux5185ux5bb9ux7269ux7406ux5386ux53f2ux4e0eux65f6ux95f4ux6743ux91cdux7684ux89e3ux91ca}

\Needspace{6\baselineskip}
\subsubsection{E10.1
学习后的参考可以包含时间组织}\label{e10.1-ux5b66ux4e60ux540eux7684ux53c2ux8003ux53efux4ee5ux5305ux542bux65f6ux95f4ux7ec4ux7ec7}

Ekman 分析区分了学习后参考包含的内容与它的时间方向 {[}E3{]}。对 35
名参与者的作者 \(8\times8\) 定位概率矩阵，删除八个对角项，保留 56
个异位置单元。回归控制提示位置偏好、输出位置偏好、圆环上的四种物理距离，以及序列位置对的平均水平；另加入时间接近和方向两个对比。前者编码为学习序列距离绝对值的负数，后者对未来位置取
\(+1\)、过去位置取 \(-1\)，只在四个学习位置之间定义。加入二者使设计秩从
18 增至 20。

主要海马输入是原文 Figure 6 notebook 实际调用的 HPa 数组。HPa 与 Figure
5 使用的 HP
是不同派生输出，不能仅凭文件名指定其解剖亚区。以参与者为单位进行 20,000
次 bootstrap，并以 100,000 次符号翻转计算检验。海马、V1
的两项对比构成四项主要比较，采用 Holm 校正。

\Needspace{9\baselineskip}

\begingroup
\small

\begin{longtable}[]{@{}
  >{\raggedright\arraybackslash}p{(\columnwidth - 6\tabcolsep) * \real{0.2000}}
  >{\raggedleft\arraybackslash}p{(\columnwidth - 6\tabcolsep) * \real{0.2667}}
  >{\raggedleft\arraybackslash}p{(\columnwidth - 6\tabcolsep) * \real{0.2667}}
  >{\raggedleft\arraybackslash}p{(\columnwidth - 6\tabcolsep) * \real{0.2667}}@{}}
\toprule\noalign{}
\begin{minipage}[b]{\linewidth}\raggedright
主要来源与对比
\end{minipage} & \begin{minipage}[b]{\linewidth}\raggedleft
平均系数
\end{minipage} & \begin{minipage}[b]{\linewidth}\raggedleft
参与者 bootstrap 95\% 区间
\end{minipage} & \begin{minipage}[b]{\linewidth}\raggedleft
Holm 校正符号翻转 \(p\)
\end{minipage} \\
\midrule\noalign{}
\endhead
\bottomrule\noalign{}
\endlastfoot
海马 Figure 6 数组：时间接近 & .0126685 & {[}.0049010, .0206462{]} &
.0164798 \\
海马 Figure 6 数组：方向 & \(-.0004862\) & {[}\(-.0027173\), .0016285{]}
& 1 \\
V1：时间接近 & \(-.0017576\) & {[}\(-.0052556\), .0018532{]} & 1 \\
V1：方向 & \(-.0002642\) & {[}\(-.0016733\), .0011401{]} & 1 \\
\end{longtable}

\endgroup

时间接近系数的单位为每个序列步对应的概率单位；方向系数为所指定带符号对比上的概率单位。结果支持该海马参考中的时间组织，方向成分仍未确定。重抽样单位是参与者，56
个单元是一个人内部的观测。来源表为
\path{historical_localizer_unique_group.csv}。

\Needspace{6\baselineskip}
\subsubsection{E10.2
实际时间表与反应重叠}\label{e10.2-ux5b9eux9645ux65f6ux95f4ux8868ux4e0eux53cdux5e94ux91cdux53e0}

对有连续 fMRI
的两位参与者，实际事件时刻在指定血流动力学响应函数下确定一个线性混合矩阵。响应均值为
6 s、标准差为 3 s，采用未滤波线性时不变计算和 3–13.5 s
提取窗时，最大非对角贡献相对平均对角贡献为 S01 的 20.75\% 和 S02 的
20.98\%。相应中心化七维混合矩阵仍具有完整对比秩，条件数为 1.298 和
1.405。这些是时间表与假定响应核共同确定的性质，不是对神经活动中由前一刺激引起比例的直接估计。

预测评价提出另外一个问题。在真实 V1
测量的留出轮次中，当前身份之外加入前一刺激身份，使汇总 SSE 在 S01 改变
\(-3.283\)\%、S02 改变
+.616\%；正数表示改善。一个人的八折共享训练资料。因此，时间表显示的潜在重叠，不会直接确定历史模型对实际数据的预测收益。不同响应核、滤波及提取窗的结果保存在
\path{historical_actual_history_mixing_summary.csv} 和
\path{historical_localizer_history_prediction_summary.csv}。

\Needspace{6\baselineskip}
\subsubsection{E10.3
放电测量中的细时间历史}\label{e10.3-ux653eux7535ux6d4bux91cfux4e2dux7684ux7ec6ux65f6ux95f4ux5386ux53f2}

Allen 分析使用一只小鼠、一次 Visual Coding Neuropixels 记录，包括 118
张自然图像和空白条件。固定 VISp 子集包含 110 个按质量标记选择的单位。从
5,950
次呈现中，删除无效时间段、空白目标，以及时间组边界和刺激中断前后三个事件的保护范围，最终保留
4,559 个目标事件。四个测试组分别含 667、1,463、1,466、963
个事件；空白仍保留在历史预测变量中。每次外层拟合用三个时间组训练，岭惩罚与历史特征缩放均在这些训练组内部选择。

模型分别以当前图像身份（C）、当前与前一图像身份（H1）、当前与前三个图像身份（H3），或当前与下一图像身份（F1）预测放电计数。历史增益允许为零，所以每个扩展模型均包含
C。报告量是相对 C 的标准化测试 SSE 减少比例。放电计数按互不重叠的 25 ms
窗提取；下一刺激最小实际起始间隔为 250.188498 ms，因此 0–250 ms
的全部窗口均位于其出现之前。

\Needspace{9\baselineskip}

\begingroup
\small

\begin{longtable}[]{@{}lrr@{}}
\toprule\noalign{}
刺激后时窗 & H1 相对 C 的 SSE 减少 & F1 相对 C 的 SSE 减少 \\
\midrule\noalign{}
\endhead
\bottomrule\noalign{}
\endlastfoot
0–25 ms & 1.1156\% & \(-.0024\)\% \\
25–50 ms & 1.3945\% & \(-.0031\)\% \\
50–75 ms & 1.0083\% & .0000\% \\
75–100 ms & .1644\% & \(-.0029\)\% \\
\end{longtable}

\endgroup

前三个时窗中，H1
在全部四个外层折中均改善预测。这些小增量表明，在此记录中，控制当前身份后仍有可预测的历史信息；残余感觉反应、适应和其他历史依赖过程都可能参与。四折刻画的是同一动物的记录。完整时窗与逐折结果保存在
\path{historical_allen_timebins_summary.csv} 和
\path{historical_allen_timebins_fold_improvements.csv}。

\Needspace{6\baselineskip}
\subsubsection{E10.4
将实测反应核迁移到快速序列}\label{e10.4-ux5c06ux5b9eux6d4bux53cdux5e94ux6838ux8fc1ux79fbux5230ux5febux901fux5e8fux5217}

FASTIMAGES 连续 MEG 示例使用 sub-01、102 个磁强计、100 Hz
采样，以及八个交替的定位/序列块。定位器提供 160 次普通刺激，40
次干扰刺激不用于核估计。快速序列共 64 个五图序列，名义刺激起始间隔为
132、164、228、612 ms；每次预测使用实际起始时间。

每次留出一个完整块，用其余七块分别估计每个刺激类别、每个传感器的 0–1 s
响应。通道尺度由训练定位器残差确定；每个序列采用自身刺激前基线。一种候选在下一物理刺激出现时截断当前刺激响应，另一种把全部已经出现刺激的完整估计响应叠加。两者共享随序列位置变化的非负增益，以及由训练资料估计的共同传感器偏置。

\Needspace{10\baselineskip}

\begingroup
\small

\begin{longtable}[]{@{}
  >{\raggedright\arraybackslash}p{(\columnwidth - 4\tabcolsep) * \real{0.2727}}
  >{\raggedleft\arraybackslash}p{(\columnwidth - 4\tabcolsep) * \real{0.3636}}
  >{\raggedleft\arraybackslash}p{(\columnwidth - 4\tabcolsep) * \real{0.3636}}@{}}
\toprule\noalign{}
\begin{minipage}[b]{\linewidth}\raggedright
刺激起始间隔
\end{minipage} & \begin{minipage}[b]{\linewidth}\raggedleft
完整响应叠加：留出 \(R^2\)
\end{minipage} & \begin{minipage}[b]{\linewidth}\raggedleft
下一刺激开始即截断：留出 \(R^2\)
\end{minipage} \\
\midrule\noalign{}
\endhead
\bottomrule\noalign{}
\endlastfoot
132 ms & .031142 & .014922 \\
164 ms & .013198 & .020630 \\
228 ms & .026549 & .041725 \\
612 ms & .003114 & .009622 \\
全部序列 & .015091 & .019519 \\
\end{longtable}

\endgroup

较丰富的反应重叠模型在最短间隔改善预测，在其余间隔未改善。在完整感觉历史基线之外加入非负前向、后向或对称激活模板，汇总
SSE 分别改变
\(-.00295\)\%、\(-.03332\)\%、\(-.01346\)\%。这些模板使用各五图序列的物理次序，折扣网格为
.1–.9，并只在训练中选择。它们校准的是这些测量中的指定模型比较；物理序列本身不提供独立确立的已学习认知图。

在固定折扣和单位预测强度下，投影去除指定历史干扰后，八折平均保留
88.94–97.13\%
的前后向模板距离。该几何保留程度依赖当前设计和固定参数比较，不代表整个连续族已被排除，也不代表实际噪声下的统计功效。定位器反应核向新任务迁移本身也需要预测验证。来源为
\path{historical_fastimages_kernel_summary.csv} 和
\path{historical_fastimages_kernel_discrimination_geometry.csv}。

\Needspace{6\baselineskip}
\subsubsection{E10.5
循环图的精确例子：有限传播坐标不等于真实时滞权重}\label{e10.5-ux5faaux73afux56feux7684ux7cbeux786eux4f8bux5b50ux6709ux9650ux4f20ux64adux5750ux6807ux4e0dux7b49ux4e8eux771fux5b9eux65f6ux6edeux6743ux91cd}

一个独立的代数例子说明，使用时间矩检验前还需要什么条件。对于真正的正折扣混合，物理时滞权重满足
\(w_1w_3-w_2^2\geq0\)
等矩不等式。有限传播基中估计出来的坐标却未必就是这些
\(w_h\)，因为更长时滞可以重新回到同一种传播形状。

\par
\noindent\begin{minipage}{\linewidth}

令 \(P\) 为 \(n\geq5\) 个状态的循环，\(A=CPC^\top\) 作用于完整中心化的
\((n-1)\) 维对比空间，则

\[
A^n=I,
\qquad A^{n-1}=-\sum_{k=0}^{n-2}A^k.
\]

\end{minipage}\par

\par
\noindent\begin{minipage}{\linewidth}

将无限等比级数按模 \(n\) 的余数分组，再消去 \(A^{n-1}\)，得到

\[
(I-\gamma A)^{-1}
=\sum_{k=0}^{n-2}b_kA^k,
\qquad
b_k=\frac{\gamma^k-\gamma^{n-1}}{1-\gamma^n}.
\]

\end{minipage}\par

\par
\noindent\begin{minipage}{\linewidth}

有限基线性独立，因此取参考 \(B=I\)，可以唯一识别所有有限系数。然而，

\[
b_1b_3-b_2^2
=-\frac{\gamma^n(1-\gamma)^2}{(1-\gamma^n)^2}<0,
\qquad 0<\gamma<1.
\]

\end{minipage}\par

当 \(n=5\)、\(\gamma=1/2\)，该值精确为
\(-8/961\)。消去当前方向后，\(AB,A^2B,A^3B\) 构成的残差设计秩为
3，因此三个有限坐标确实可以唯一估计，但矩行列式仍为负。真实无限核系数是
\(w_1=\gamma\)、\(w_2=\gamma^2\)、\(w_3=\gamma^3\)，满足
\(w_1w_3-w_2^2=0\)。有限坐标的负行列式完全由尾部折叠进有限基产生。

要把拟合系数解释成实际时滞权重，测量设计就必须把这些时滞与尾部分开。若瞬态算子满足
\(N^4=0\)，且有限设计可辨识，则更高传播幂全部消失，其三个正时滞系数可以直接解释。独立任务操纵，或显式且可辨识的尾部模型，也能提供这样的区分。上述一般循环反例同样适用于校准中的
12 状态和 40 状态循环；E9 的 resolvent
拟合直接建模完整级数，没有对有限坐标施加该不适用的矩检验。完整归档推导见
\path{reference_archive/cyclic_lag_scope_proof.md}。

\Needspace{7\baselineskip}
\subsection{E11.
研究覆盖与来源对应}\label{e11.-ux7814ux7a76ux8986ux76d6ux4e0eux6765ux6e90ux5bf9ux5e94}

研究计划考察了互补的测量情境。下表记录每个分支的实际可用材料及其在论证中的作用。材料不足的分支保留为数据资格记录。各分支的样本量不能合并：Ekman
的两位连续 fMRI 参与者属于其 35 人队列，FASTIMAGES 的 sub-01 属于其 30
人 MEG 队列，Huang 的 EEG 与 fMRI 是同一批 33 名参与者。

\begingroup
\small

\begin{longtable}[]{@{}
  >{\raggedright\arraybackslash}p{(\columnwidth - 6\tabcolsep) * \real{0.2500}}
  >{\raggedright\arraybackslash}p{(\columnwidth - 6\tabcolsep) * \real{0.2500}}
  >{\raggedright\arraybackslash}p{(\columnwidth - 6\tabcolsep) * \real{0.2500}}
  >{\raggedright\arraybackslash}p{(\columnwidth - 6\tabcolsep) * \real{0.2500}}@{}}
\toprule\noalign{}
\begin{minipage}[b]{\linewidth}\raggedright
分支
\end{minipage} & \begin{minipage}[b]{\linewidth}\raggedright
实际材料与分析单位
\end{minipage} & \begin{minipage}[b]{\linewidth}\raggedright
对本文论证的作用
\end{minipage} & \begin{minipage}[b]{\linewidth}\raggedright
原始来源
\end{minipage} \\
\midrule\noalign{}
\endhead
\bottomrule\noalign{}
\endlastfoot
D01，Tarder-Stoll & 同一人的四次定位、两次预测扫描；任务与归档字段审查 &
确定任务专属状态矩阵的资格；现有事件/路径映射不足以重建路径专属算子 &
\href{https://doi.org/10.1038/s41467-024-53293-3}{论文}；OpenNeuro
ds005125 \\
D02，Ekman & 35 人作者派生数组；其中两人的连续 fMRI &
学习后参考内容、整提示留出比较、参考预测和物理历史校准；E7、E10.1–E10.2
&
\href{https://doi.org/10.7554/eLife.78904}{论文}；\href{https://doi.org/10.34973/bsy6-9h29}{数据} \\
D03，Kern & DSMR101 的 MEG 定位器与目标测量 &
参考—目标别名构造、连续族校准和冻结读出迁移；E6、E9 &
\href{https://doi.org/10.7554/eLife.93357}{论文}；\href{https://doi.org/10.5281/zenodo.8001755}{数据} \\
D04，He & 26 名可分析参与者的 EEG 解码/RDM/RSA；30 人供应集作敏感性分析
& 在现有概率/几何输出层面检验学习相关内容和直接方向对比 &
\href{https://doi.org/10.1371/journal.pbio.3003740}{论文}；\href{https://osf.io/sn9kv/}{OSF} \\
D05，THINGS-fMRI & 一人、100 张图像、12 次记录、72 个预选体素 &
全空间与低维预测、已知变换校准、分离记录的阈值评价 &
\href{https://github.com/ViCCo-Group/THINGS-data}{THINGS-data}；OpenNeuro
ds004192 \\
D06，THINGS-MEG & 一人两次记录各一个 run；五张共同图像，各保留一次呈现 &
跨测量对应关系，以及现有重复结构的资格检查 &
\href{https://github.com/ViCCo-Group/THINGS-data}{THINGS-data} \\
D07，CNeuroMod THINGS & 一人、36 次记录、11,646 个保留试次、七个 ROI &
在供应体素选择条件下进行跨图像、跨记录通道预测 &
\href{https://github.com/courtois-neuromod/cneuromod-things}{官方代码}；\href{https://docs.cneuromod.ca/datasets/things.html}{数据说明} \\
D08，Turner & ID13 两次 EEG 定位记录，各 1,980 个合格事件 &
时间分离的参考预测、任务相容读出、实际块限制和迁移；E2–E6、E9 &
\href{https://doi.org/10.1371/journal.pbio.3003189}{论文} \\
D09，FASTIMAGES-MEG & 30 人概率输出；sub-01 连续 MEG &
物理序列时序、实测响应核迁移和历史模板比较；E10.4 &
\href{https://github.com/skjerns/FASTIMAGES}{官方数据与代码} \\
D10，FASTIMAGES-fMRI & 40 人概率输出、3,000 个试次，每试次 13 TR &
概率输出层面的时间校准 &
\href{https://github.com/skjerns/FASTIMAGES}{官方数据与代码} \\
D11，缺逐试次 MEG 的分支 & 审查材料中没有合格公开逐试次 MEG/定位器文件 &
数据可用性边界；未实施经验检验 & 来源状态保留于研究报告 \\
D12，Kahn & 322 人、535 段完整图学习序列 &
行为读出竞争、近期性与局部痕迹替代解释、跨人及跨图预测 &
\href{https://doi.org/10.1371/journal.pcbi.1013696}{论文}；\href{https://osf.io/jupbz/}{OSF} \\
D13，Garvert & 23 人 ROI 配对反应；26 人处理后反应时 &
非均匀无向图的平稳时间逆过程、候选几何与整 run 留出比较 &
\href{https://doi.org/10.7554/eLife.17086}{论文} \\
D14，Allen Neuropixels & 一鼠、一次记录、5,950
次自然场景呈现；细时间窗用 110 个 VISp 单位 &
实测感觉历史、测量时窗和秩/噪声校准；E10.3 & Allen Visual Coding
Neuropixels；文件 \path{sub-707296975_ses-721123822.nwb} \\
D15，源档案不可访问分支 & 审查时原 Dryad 档案不可用，未取得等价输入 &
数据可用性边界；未实施经验检验 & 来源状态保留于研究报告 \\
D16，Huang & 33 人 EEG 输出及同一批 33 人 fMRI 输出；2,561 个 fMRI 试次
& 概率归一化约束、试次边界核对与序列统计 &
\href{https://doi.org/10.1038/s41467-024-51582-5}{论文}；\href{https://zenodo.org/records/12547774}{数据} \\
\end{longtable}

\endgroup

正式第一轮报告 \path{SR_real_data_full_report_20261123.pdf}
及其主张—证据表记录全部 16
个分支。第二轮报告和附录保留参考、时间历史、设计秩及已知变换分析。连续族归档
\path{SR_dFB_complete_20261124.zip} 提供 E9 使用的逐折 dFB
记录、独立复原与连续优化核对。这些报告及紧凑来源表保存在
\path{reference_archive/} 和
\path{source_data/}；原始归档路径与哈希记录于经验来源清单。当前事件数组校准另有
\path{source_data/README_fig07.md} 说明。

这份目录保留完整研究覆盖，包括停留在测量资格阶段的分析，并使每一项经验主张对应到自己的数据和统计单位。原始论文负责数据身份与实验设计的来源归属；本文引入的数值结果属于上述归档再分析。

\Needspace{6\baselineskip}
\subsubsection{经验资料来源}\label{ux7ecfux9a8cux8d44ux6599ux6765ux6e90}

{[}E1{]} Turner W, Sexton C, Johnson PA, Wilson EM, Hogendoorn H.
Predictable motion is progressively extrapolated across temporally
distinct processing stages in the human visual cortex. \emph{PLOS
Biology} 23, e3003189 (2025).
\url{https://doi.org/10.1371/journal.pbio.3003189}

{[}E2{]} Kern S. Reactivation strength during cued recall is modulated
by graph distance within cognitive maps. Public data archive, version 1
(2023). \url{https://doi.org/10.5281/zenodo.8001755}

{[}E3{]} Ekman M, Kusch S, de Lange FP. Successor-like representation
guides the prediction of future events in human visual cortex and
hippocampus. \emph{eLife} 12, e78904 (2023).
\url{https://doi.org/10.7554/eLife.78904}

{[}E4{]} Kern S, Nagel J, Gerchen MF, Gürsoy Ç, Meyer-Lindenberg A,
Kirsch P, Dolan RJ, Gais S, Feld GB. Reactivation strength during cued
recall is modulated by graph distance within cognitive maps.
\emph{eLife} 12, RP93357 (2024).
\url{https://doi.org/10.7554/eLife.93357}

\clearpage

\section{C.
与相关研究的关系}\label{c.-ux4e0eux76f8ux5173ux7814ux7a76ux7684ux5173ux7cfb}

\Needspace{7\baselineskip}
\subsection{C1.
预测地图、时间背景与方向}\label{c1.-ux9884ux6d4bux5730ux56feux65f6ux95f4ux80ccux666fux4e0eux65b9ux5411}

Dayan 提出后继表征，以折扣占据预测将环境转移结构与奖励分开
{[}1{]}。Stachenfeld
等进一步发展了这一结构对海马群体编码、位置场和表征相似性的解释
{[}2{]}。本文从这种机制解释出发，追问其中哪些性质能够经过观测过程而被保留。

时间方向也已通过学习规则与行为实验得到研究。Gershman 等联系了 SR
学习与时间背景模型，并分析了二者何时给出共同预测、何时分离 {[}3{]}。Keck
等研究局部学习规则的对称性如何改变预测表征与泛化 {[}10{]}。Momennejad
等利用重估区分规划策略 {[}4{]}；Sharp 和 Eldar
操纵前向与后向预测支持不同决策的实验条件
{[}9{]}。这些工作确立了竞争机制与实验干预的价值。本文的互反任务结果进一步给出未知神经特征映射下方向性质的充分条件，并将跨条件切换机制的可能性纳入竞争解释。混合地图与混合转移之间的区别由
T22 明确给出。

\Needspace{7\baselineskip}
\subsection{C2.
神经观测与参考测量}\label{c2.-ux795eux7ecfux89c2ux6d4bux4e0eux53c2ux8003ux6d4bux91cf}

关系几何、未来相关内容和时间尺度是相互联系的经验目标。Garvert
等研究海马—内嗅系统中的抽象关系知识 {[}5{]}。Brunec 和 Momennejad
考察海马与前额叶层级中的预测表征 {[}6{]}。Ekman
等结合序列学习、省略探测和独立定位测量比较未来与过去内容
{[}7{]}。Tarder-Stoll 等研究多步预期中过去与未来的层级表征 {[}8{]}；He
等研究学习后睡眠与后继类别证据之间的关系 {[}19{]}。

本文把这些观测层次连接到具体数学目标：全部相容解释共享的方向、非负混合的净方向、折扣范围，或有限的时间权重序列。已学习参考会将自身的预测结构带入两次测量之间的算子。测量的时间持续性还提供了产生表观后向结构的另一种机制；T23
给出其均值关系以及独立探测的作用。

经验来源保留各自原有的科学目标。Turner 的 EEG
研究考察不同加工阶段的运动外推 {[}18{]}，Kern 的 MEG
研究考察线索回忆中的图距离
{[}25{]}。这些来源的重复测量为本文提供参考估计与已知变换校准的基础。

\Needspace{7\baselineskip}
\subsection{C3.
表征模型与统计辨识}\label{c3.-ux8868ux5f81ux6a21ux578bux4e0eux7edfux8ba1ux8fa8ux8bc6}

Diedrichsen 和 Kriegeskorte 将编码、模式成分和表征相似性方法纳入共同框架
{[}12{]}。Gardner 和 Liu，以及 Sprague
等，解释了编码基底与模型假设如何影响重建响应的解释
{[}13,14{]}。Kriegeskorte 等分析了循环选择问题 {[}15{]}；Cai
等考察任务结构对表征相似性的偏差及其贝叶斯处理 {[}16{]}。

本文的任务相容压缩保留任务的指定作用，而不要求恢复全部神经特征。相容集合定义在带标签的完整模型族上，因此相同数值预测仍可保留不同机制标签。性质辨识由此考察所有保留解释是否一致。在
T26–T27 的抽样条件下，Hotelling 的多变量均值分布 {[}26{]}
提供有限样本的候选残差阈值。这一步将总体关系转为具有条件错误控制的声明；实际重复单元和协方差继续作为实验设计的一部分。

\Needspace{7\baselineskip}
\subsection{C4.
时间恢复与更广的建模纲领}\label{c4.-ux65f6ux95f4ux6062ux590dux4e0eux66f4ux5e7fux7684ux5efaux6a21ux7eb2ux9886}

Momennejad 和 Howard 从已知尺度的多尺度 SR
出发，通过近似逆拉普拉斯变换恢复粗略的未来时间线
{[}11{]}。本文的定量结果处理未知折扣和权重。独立传播形状或不同的可见旋转块约束共享参数；基于二元
de Bruijn 构造 {[}27{]}
的有限记忆任务则将可观测滞后分配到不同矩阵位置。直接参考和已学习反向参考进一步允许不同形式的核恢复。

Discussion 将这些结果联系到成功拟合究竟确立了什么这一一般问题
{[}17,20{]}。其中的时间知觉例子来自 Youguo Chen 等的模型
{[}21,22{]}；连续参数补偿则由本文的 SR
族给出解析实例。李群对称的研究方向与既有基于对称性的可辨识性和参数化工作相联系
{[}23,24{]}。本文数学结果针对明确规定的 SR
族，对称性纲领为进一步构造和比较模型类提供研究方向。

\clearpage

\section{D. 参考文献}\label{d.-ux53c2ux8003ux6587ux732e}

\fontsize{9}{11}\selectfont
\setlength{\parskip}{2pt}
\setlist[enumerate]{itemsep=2pt}

\begin{enumerate}
\def\labelenumi{\arabic{enumi}.}
\item
  Dayan, P. Improving Generalization for Temporal Difference Learning:
  The Successor Representation. \emph{Neural Computation} \textbf{5},
  613–624 (1993). \url{https://doi.org/10.1162/neco.1993.5.4.613}
\item
  Stachenfeld, K. L., Botvinick, M. M. \& Gershman, S. J. The
  hippocampus as a predictive map. \emph{Nature Neuroscience}
  \textbf{20}, 1643–1653 (2017). Corrected version; correction:
  10.1038/s41593-018-0133-1. \url{https://doi.org/10.1038/nn.4650}
\item
  Gershman, S. J., Moore, C. D., Todd, M. T., Norman, K. A. \&
  Sederberg, P. B. The Successor Representation and Temporal Context.
  \emph{Neural Computation} \textbf{24}, 1553–1568 (2012).
  \url{https://doi.org/10.1162/NECO_a_00282}
\item
  Momennejad, I., Russek, E. M., Cheong, J. H., Botvinick, M. M., Daw,
  N. D. \& Gershman, S. J. The successor representation in human
  reinforcement learning. \emph{Nature Human Behaviour} \textbf{1},
  680–692 (2017). \url{https://doi.org/10.1038/s41562-017-0180-8}
\item
  Garvert, M. M., Dolan, R. J. \& Behrens, T. E. J. A map of abstract
  relational knowledge in the human hippocampal–entorhinal cortex.
  \emph{eLife} \textbf{6}, e17086 (2017).
  \url{https://doi.org/10.7554/eLife.17086}
\item
  Brunec, I. K. \& Momennejad, I. Predictive representations in
  hippocampal and prefrontal hierarchies. \emph{Journal of Neuroscience}
  \textbf{42}, 299–312 (2022).
  \url{https://doi.org/10.1523/JNEUROSCI.1327-21.2021}
\item
  Ekman, M., Kusch, S. \& de Lange, F. P. Successor-like representation
  guides the prediction of future events in human visual cortex and
  hippocampus. \emph{eLife} \textbf{12}, e78904 (2023).
  \url{https://doi.org/10.7554/eLife.78904}
\item
  Tarder-Stoll, H., Baldassano, C. \& Aly, M. The brain hierarchically
  represents the past and future during multistep anticipation.
  \emph{Nature Communications} \textbf{15}, 9094 (2024).
  \url{https://doi.org/10.1038/s41467-024-53293-3}
\item
  Sharp, P. B. \& Eldar, E. Humans adaptively deploy forward and
  backward prediction. \emph{Nature Human Behaviour} \textbf{8},
  1726–1737 (2024). Publisher correction: 10.1038/s41562-024-01978-6.
  \url{https://doi.org/10.1038/s41562-024-01930-8}
\item
  Keck, J., Barry, C., Doeller, C. F. \& Jost, J. Impact of symmetry in
  local learning rules on predictive neural representations and
  generalization in spatial navigation. \emph{PLOS Computational
  Biology} \textbf{21}, e1013056 (2025).
  \url{https://doi.org/10.1371/journal.pcbi.1013056}
\item
  Momennejad, I. \& Howard, M. W. Predicting the Future with Multi-scale
  Successor Representations. \emph{bioRxiv} (2018). Preprint, version 1.
  \url{https://doi.org/10.1101/449470}
\item
  Diedrichsen, J. \& Kriegeskorte, N. Representational models: A common
  framework for understanding encoding, pattern-component, and
  representational-similarity analysis. \emph{PLOS Computational
  Biology} \textbf{13}, e1005508 (2017).
  \url{https://doi.org/10.1371/journal.pcbi.1005508}
\item
  Gardner, J. L. \& Liu, T. Inverted Encoding Models Reconstruct an
  Arbitrary Model Response, Not the Stimulus. \emph{eNeuro} \textbf{6},
  ENEURO.0363-18.2019 (2019).
  \url{https://doi.org/10.1523/ENEURO.0363-18.2019}
\item
  Sprague, T. C., Boynton, G. M. \& Serences, J. T. The Importance of
  Considering Model Choices When Interpreting Results in Computational
  Neuroimaging. \emph{eNeuro} \textbf{6}, ENEURO.0196-19.2019 (2019).
  \url{https://doi.org/10.1523/ENEURO.0196-19.2019}
\item
  Kriegeskorte, N., Simmons, W. K., Bellgowan, P. S. F. \& Baker, C. I.
  Circular analysis in systems neuroscience—the dangers of double
  dipping. \emph{Nature Neuroscience} \textbf{12}, 535–540 (2009).
  \url{https://doi.org/10.1038/nn.2303}
\item
  Cai, M. B., Schuck, N. W., Pillow, J. W. \& Niv, Y. Representational
  structure or task structure? Bias in neural representational
  similarity analysis and a Bayesian method for reducing bias.
  \emph{PLOS Computational Biology} \textbf{15}, e1006299 (2019).
  \url{https://doi.org/10.1371/journal.pcbi.1006299}
\item
  Wilson, R. C. \& Collins, A. G. E. Ten simple rules for the
  computational modeling of behavioral data. \emph{eLife} \textbf{8},
  e49547 (2019). \url{https://doi.org/10.7554/eLife.49547}
\item
  Turner, W., Sexton, C., Johnson, P. A., Wilson, E. M. \& Hogendoorn,
  H. Predictable motion is progressively extrapolated across temporally
  distinct processing stages in the human visual cortex. \emph{PLOS
  Biology} \textbf{23}, e3003189 (2025).
  \url{https://doi.org/10.1371/journal.pbio.3003189}
\item
  He, X., Büchel, P. K., Faghel-Soubeyrand, S., Klingspohr, J., Kehl, M.
  S. \& Staresina, B. P. Sleep strengthens successor representations of
  learned sequences in humans. \emph{PLOS Biology} \textbf{24}, e3003740
  (2026). \url{https://doi.org/10.1371/journal.pbio.3003740}
\item
  Roberts, S. \& Pashler, H. How persuasive is a good fit? A comment on
  theory testing. \emph{Psychological Review} \textbf{107}, 358–367
  (2000). \url{https://doi.org/10.1037/0033-295X.107.2.358}
\item
  Chen, Y., Peng, C. \& Avitt, A. A unifying Bayesian framework
  accounting for spatiotemporal interferences with a deceleration
  tendency. \emph{Vision Research} \textbf{187}, 66–74 (2021).
  \url{https://doi.org/10.1016/j.visres.2021.06.005}
\item
  Chen, Y., Yu, J., Hou, C. \& Peng, C. Direct evidence for logarithmic
  magnitude representation in the central nervous system.
  \emph{Neuroscience} \textbf{561}, 127–138 (2024).
  \url{https://doi.org/10.1016/j.neuroscience.2024.10.025}
\item
  Merkt, B., Timmer, J. \& Kaschek, D. Higher-order Lie symmetries in
  identifiability and predictability analysis of dynamic models.
  \emph{Physical Review E} \textbf{92}, 012920 (2015).
  \url{https://doi.org/10.1103/PhysRevE.92.012920}
\item
  Popovych, R. O. \& Bihlo, A. Symmetry preserving parameterization
  schemes. \emph{Journal of Mathematical Physics} \textbf{53}, 073102
  (2012). \url{https://doi.org/10.1063/1.4734344}
\item
  Kern, S., Nagel, J., Gerchen, M. F., Gürsoy, Ç., Meyer-Lindenberg, A.,
  Kirsch, P., Dolan, R. J., Gais, S. \& Feld, G. B. Reactivation
  strength during cued recall is modulated by graph distance within
  cognitive maps. \emph{eLife} \textbf{12}, RP93357 (2024).
  \url{https://doi.org/10.7554/eLife.93357}
\item
  Hotelling, H. The generalization of Student's ratio. \emph{Annals of
  Mathematical Statistics} \textbf{2}, 360–378 (1931).
  \url{https://doi.org/10.1214/aoms/1177732979}
\item
  de Bruijn, N. G. A combinatorial problem. \emph{Proceedings of the
  Section of Sciences of the Koninklijke Nederlandse Akademie van
  Wetenschappen te Amsterdam} \textbf{49}, 758–764 (1946).
  \url{https://research.tue.nl/en/publications/a-combinatorial-problem/}
\end{enumerate}
\end{document}
